\documentclass[11pt,letterpaper]{article}
\usepackage[T1]{fontenc}
\usepackage[utf8]{inputenc}
\usepackage{lmodern,amsmath,amssymb,amsthm,mathtools}
\usepackage[margin=0.87in]{geometry}
\usepackage{microtype,booktabs,longtable,array,enumitem,xcolor,fancyhdr,xurl}
\usepackage[colorlinks=true,linkcolor=blue!40!black,urlcolor=blue!45!black,citecolor=blue!40!black]{hyperref}
\usepackage{listings}
\hypersetup{pdftitle={PA194AX11QH: a new 194-state proof-search machine},pdfauthor={Research report prepared for Vladimir Reshetnikov}}
\setlist{nosep,leftmargin=*}
\setlength{\parskip}{4pt}
\setlength{\parindent}{0pt}
\setlength{\emergencystretch}{3em}
\pagestyle{fancy}\fancyhf{}\fancyhead[L]{\small PA194AX11QH}\fancyhead[R]{\small 14 November 2026}\fancyfoot[C]{\thepage}
\setlength{\headheight}{14pt}
\lstset{basicstyle=\ttfamily\small,breaklines=true,columns=fullflexible,keepspaces=true,frame=single,rulecolor=\color{black!20}}
\newtheorem{theorem}{Theorem}[section]
\newtheorem{lemma}[theorem]{Lemma}
\newtheorem{proposition}[theorem]{Proposition}
\newtheorem{corollary}[theorem]{Corollary}
\theoremstyle{definition}\newtheorem{definition}[theorem]{Definition}
\theoremstyle{remark}\newtheorem{remark}[theorem]{Remark}
\newcommand{\PA}{\mathrm{PA}}
\newcommand{\iPRA}{\mathrm{iPRA}}
\newcommand{\PRA}{\mathrm{PRA}}
\newcommand{\Prf}{\operatorname{Prf}}
\newcommand{\Con}{\operatorname{Con}}
\newcommand{\NH}{\operatorname{NH}}
\newcommand{\Halt}{\operatorname{Halt}}
\newcommand{\file}[1]{\path{#1}}
\newcommand{\MX}{M_{\mathrm{AX11QH}}}
\newcommand{\Old}{M_{194}}
\newcommand{\snapshot}{\path{1c9524d8bf2e0a27b9a16ff622359927f536e663}}
\title{PA194AX11QH:\\A new 194-state proof-search machine\\[5pt]\large Corpus comparison, literal-tape certificates,\\and the remaining iPRA consistency calibration}
\author{Research report prepared for Vladimir Reshetnikov}
\date{14 November 2026}
\begin{document}
\maketitle
\begin{abstract}
We compare the current ProveIt PA-consistency-machine corpus at the pinned
commit \snapshot. The 72 deliveries contain 82 distinct complete spanning
194-state transition graphs, including 40 principal construction claims;
names alone substantially overcount or misidentify designs. We construct
\textbf{PA194AX11QH}, a new 194-working-state binary single-tape machine
combining an eleven-entry axiom-command protocol, exchanged universal-
distribution premises, no-push modus ponens, and a cached-only contradiction
test. Exact source reconstruction, 28 simultaneous fold cases, all 4,778
arbitrary-register source outcomes, four list-sweep lemmas derived directly
from the rows, and complete 57,922-step blank initialization pass independent
executable checks. Exact graph comparison finds no match in the current
frontier, even under broad structural symmetries.

\textbf{The original target remains incomplete.} No blank-start machine
below 194 states is established. More significantly, the supplied physical
PA194 machines do not already have the claimed iPRA calibration in the
pinned repository: its README explicitly leaves both machine-specific
iPRA directions open. This report separates the new concrete construction,
relative preservation arguments, and finite certificates from that missing
PA-refutation calibration. It does not present a conditional theorem or a
successful Python run as a completed proof of
$\iPRA\vdash\NH(\MX)\leftrightarrow\Con(\PA)$.
\end{abstract}
\tableofcontents
\clearpage

\section{Result and exact conventions}

The requested target is a literal Turing machine $M$ with fewer than 194
working states, or a materially different 194-state construction, together
with a proof in \emph{intuitionistic} primitive recursive arithmetic of
\[
 \NH(M)\quad\longleftrightarrow\quad\Con(\PA).
\]
The theory whose refutations are sought is classical PA. The intuitionistic
qualification concerns the base theory in which the equivalence must be
proved. These two roles cannot be exchanged.

The delivered new machine has alphabet $\{0,1\}$, a two-way infinite tape,
an all-zero initial tape, head initially at coordinate zero, and initial
state zero. Its working states are $0,\ldots,193$; the separate terminal
label is 194. Every working-state/symbol pair has exactly one transition,
and every transition moves exactly one cell left or right. No prescribed
nonblank input, auxiliary tape, stay instruction, or free initialization
is used. All 194 working states are reachable in the transition graph
from the designated initial state. Graph reachability does not assert that
every state is encountered in the actual blank run.

The normative dense table is \file{candidate/machine/machine.tm}, with
SHA-256
\begin{center}\small\path{19e215df20429fa6a35f0a99ba0dcfcb40893a6e4b9a4ae6e450b0225f110e50}\end{center}
The native table uses start 74 and terminal 254, with explicit bijection
\file{candidate/machine/labels.json}. Its SHA-256 is
\begin{center}\small\path{424bd5175bbca161e76af485f1eda03f2f169d90afe80563035b60289361bf07}\end{center}
Appendix~\ref{app:table} prints all transitions.

\begin{center}
\begin{tabular}{@{}lr@{}}\toprule
New construction and fresh checks & Result\\\midrule
Working states / binary rows & 194 / 388\\
Fixed controls + readers $-$ folds & $134+67-7$\\
Unfolded parent controls & 201\\
Simultaneous local fold cases & 28\\
Primitive source sites & 3,671\\
Source outcomes, including both decrement branches & 4,778\\
Independent symbolic execution operations & 1,130,254\\
Generic list lemmas rederived from literal rows & 4\\
Minimum positive duration constant & 3\\
Complete blank initialization & 57,922 steps\\
Existing spanning 194-state graphs compared & 82\\
Matches, including broad symmetries & 0\\\bottomrule
\end{tabular}
\end{center}

The new construction is a composition of previously proposed mechanisms.
Its contribution is the new recurring command interface, its concrete
source/table binding, corrected proof-word operations for that combination,
and independent verification. It is not a newly discovered logical axiom
basis or a state-count record.

\section{The comparison and what its counts mean}
\label{sec:comparison}

\subsection{Pinned scope and reproducibility}

The comparison is pinned to ProveIt commit \snapshot\ \cite{repo}.
There are 61 filed packets under
\file{Computability/BusyBeaver/Research/PA/docs} and 11 incoming ZIP archives.
The packet ledger records every one of these 72 deliveries, including
wrapper names, source locators, hashes, reported contributions, and the
boundary of the new audit. This is needed because several archives have
been unpacked and deduplicated in the repository while newer archives
remain incoming \cite{packets}.

The independent scanner reads all currently tracked \file{.tm},
\file{.txt}, \file{.tsv}, \file{.csv}, \file{.json}, and \file{.zip}
files in the specified PA corpus and incoming archives, including nested
archives. No file-size cutoff remains for those formats. It recognizes
complete numeric binary left/right transition tables and supported explicit
JSON transition objects. It does not infer a machine from arbitrary
program text, TeX-only tables, JSONL streams, or opaque binary formats;
the manuscript/README ledger separately covers named proposals.

\begin{center}
\begin{tabular}{@{}lr@{}}\toprule
Census quantity & Count\\\midrule
Deliveries & 72\\
JSON / TXT payloads examined & 9,400 / 1,867\\
TSV / CSV / TM payloads examined & 388 / 209 / 3,356\\
ZIP references / unique ZIP payloads & 37 / 25\\
Recognized complete table occurrences & 3,422\\
Normalized labeled tables & 752\\
Spanning 194-state transition-graph classes & 82\\
Additional disconnected 194-state diagnostics & 2\\
Complete experimental 193-state graph classes & 53\\\bottomrule
\end{tabular}
\end{center}

These counts do not directly equal the repository's historical count of
7,130 table occurrences: that older ledger counts original archive members,
including thousands of copies now omitted by deduplication, and its
recognition scope differs \cite{machineledger}. The new scanner records 22
incomplete or nonconforming \file{.tm} payloads and excludes 35 incidental
one-state JSON fragments. It retains the corresponding exception records.
No parse failure is silently converted into an accepted consistency machine.

\subsection{Identity and novelty}

For a chosen root reaching every working state, traverse successors in the
fixed read-symbol order $0,1$, assigning new labels at first encounter.
Record for every state its two triples (written symbol, direction, next
label), using a separate terminal marker. The resulting tuple is a
complete invariant of rooted transition-graph isomorphism: equality of
tuples gives the isomorphism directly by matching traversal indices.
Taking the minimum over every spanning root removes dependence on the
initial label for this structural comparison.

A broader comparison additionally permits global left/right reflection,
global exchange of the two symbols, and omission of the write/move effects
on terminal edges. The last effects do not affect whether the machine
halts. Symbol exchange is used only as a strong structural novelty test;
it need not preserve the all-zero initial tape. The designated blank start
is independently pinned for every constructed machine.

The new comparison checks full canonical tuples, not just their SHA-256
identifiers. PA194AX11QH has zero strict or broad matches among all 82
spanning representatives. The two disconnected diagnostics cannot be
isomorphic to its spanning graph and are explicitly reported separately.
Its strict graph identifier is
\begin{center}\small\path{78245b0fa67acabdd21a9a9d7b5c2ba8975ed053a06dc408fb0042f1d5363bd7}\end{center}
and its broad identifier is
\begin{center}\small\path{5ca4022d1d388653ae5f50a56a13c2ac579b6e2b7e92229dfcef3a196364e94a}.\end{center}

\subsection{Evidence categories and shared architecture}

Of the 82 spanning graphs, 40 are reported principal constructions, seven
are supplementary constructions, five are source-validated research
reassociations, 25 are other research auxiliaries or seeds, and five are
explicitly rejected. These are attributed historical evidence categories.
This work does not reprove 82 historical arithmetic calibrations.

All 40 principal designs use essentially the same architecture: enumeration
of command words acting on a persistent numerical proof store, a unary
register-file interpreter, and a compressed binary source decoder. Their
material differences occur in formula arithmetic, emitted axiom heads,
proof-store effects, source association, register placement, and known-zero
return conventions. Appendix~\ref{app:principal} compares every principal
class, and Appendix~\ref{app:frontier} identifies every spanning graph.
The detailed CSV/JSON files retain the mechanisms and source provenance.

\begin{description}[style=nextline]
\item[Original and alternate associations.]
The primary machine uses 67 reader nodes and seven contextual folds; the
alternate association uses 68 readers and eight folds. Both yield 194
working states and use the ten-entry AX protocol.
\item[C and CS.]
Successful MP places the consequent in the cached head; CS also changes
the destructive comparison order. These should not be confused with H.
\item[N and H.]
N removes MP's dummy-zero push while returning the consequent visibly.
H combines visible no-push MP with a cached-only Shift halt test. A visible
contradiction therefore needs an explicit exposure step.
\item[Q.]
Q exchanges the two curried premises of universal distribution. Its
logical proof transformation is separate from the controller reductions.
\item[D.]
The diagonal generator uses $\Gamma(D)$ rather than $\Gamma(A)$; the
selectors must be changed accordingly. PA194D and PA194DHS use this choice.
\item[Radix 5 and 6.]
The constructor coefficient changes throughout the numerical formulas.
A small transition edit can implement this extensive semantic change, but
one must transport the calculus and recheck the runtime.
\item[Known-zero returns and register permutations.]
J, F, A, X, ZR and HAZ use selected proved-zero decrements in place of jumps,
sometimes with a physical register permutation. A source invariant does
not automatically establish a valid blank initialization.
\item[AX11.]
AX pushes its final guard as well, then enters the full MP field reader.
It consumes a genuine fourth parameter and leaves eleven prefix entries.
The old three-parameter encoding is incorrect for this interface.
\end{description}

The counts within the 40 principal classes are informative: 28 use
coefficient four, six coefficient five, and six coefficient six; 12 use Q;
10 retain dummy-push MP, 12 use cached-result MP, and 18 use visible no-push
MP; 21 use the sum halt test and 19 the cache test; 14 use the symmetric
comparator. These are protocol counts, not independently checked theorem
counts.

\subsection{Duplicate and near-miss corrections}

PA194C occurs as the same principal graph in five deliveries. The lowercase
\file{pa194ns_conpa} machine equals PA194HS, whereas the uppercase
\file{PA194NS_research} machine retains the sum test and is different.
PA194QN equals PA194NQ; the two QNS deliveries agree; PA194CSZ equals
PA194X. PA194QH5 is duplicated in its dedicated archive and the AX11 archive;
PA194QN6 is duplicated in its dedicated archive and the combined QN archive.
Native/dense renumberings do not create additional machines.

The 53 current 193-state graphs include the rejected coefficient-three
calculus, the initialized runtime-only marker-fold core, and 51 further
unverified or rejected variants. Thus an old census saying that there are
only two 193-state objects cannot be used for the current snapshot. None
is supplied with an accepted blank-start PA calibration.

The coefficient-three shortcut is especially instructive. With its old
equality offset, $E(1,0)=6=I(0,1)$, and the retained numerical calculus has
an explicit refutation. The later all-offset obstruction applies only to
its specified unchanged rule/head family. Neither observation is a general
lower bound on PA-consistency-machine size. The runtime-only core has a
different defect: valid initialized source frames are simulated, but the
blank computation leaves the intended invariant and enters an outward scan.

% Standalone report fragment. Requires amsmath, amssymb, booktabs, and hyperref.
% Its relative arithmetic theorem is developed separately; this section makes
% no claim to a completed iPRA consistency calibration.
\section{The new PA194AX11QH construction}
\label{sec:ax11qh-construction}

The concrete output of the construction work is a new 194-working-state
machine, called \textbf{PA194AX11QH}. It combines four mechanisms from the
supplied packets: the exchanged universal-distribution head (Q), the
AX command that pushes its gate and falls through the full modus-ponens
reader (AX11), no-push modus ponens (N), and a cached-only contradiction
test (H). The constructor coefficient remains four. The contribution is
this precise composition, its fresh source and tape verification, and the
reduction searches described below. Each constituent mechanism is
attributed to its supplied packet in \texttt{PROVENANCE.json}.

The machine is deterministic, uses one two-way infinite binary tape,
and moves one cell left or right on every transition. Its initial tape
is entirely zero and its initial head coordinate is zero. The dense table
has initial state 0, working states $0,\ldots,193$, and a separate terminal
label 194 with no outgoing transitions. Thus there are 388 transition rows;
the terminal label is not included in the count of 194. The native table
has initial state 74 and terminal label 254. Native labels are identifiers,
not state counts.

The completed checks concern this exact machine. They do not supply the
missing baseline proof in iPRA of the PA-consistency calibration. The
relative proof-strength discussion elsewhere in this report keeps those
obligations explicit. In particular, successful Python verification is
not being reported as a newly kernel-checked Lean or Rocq theorem.

\subsection{Numerical operations and the exchanged logical head}

Use Cantor pairing
\[
 p(a,b)=\frac{(a+b)(a+b+1)}2+b,
 \qquad [a;w]=p(a,w),
\]
with total projections $\operatorname{hd}$ and $\operatorname{tl}$. In
particular, $[0;0]=0$ and both projections of zero are zero. Write
$[a_0,\ldots,a_k;w]$ for iterated pairing with $a_0$ on top. None of the
arguments below replaces this numerical stack by a free list.
The formula constructors are
\[
 \begin{aligned}
 I(A,B)&=4p(A,B), & M(a,b)&=4p(a,b)+1,\\
 U(a,B)&=4p(B,a)+2, & E(a,b)&=4p(a,b)+3.
 \end{aligned}
\]
Zero is the truth code and one is the designated contradiction code.
The alias $I(0,0)=0$ is retained throughout the command calculations.

The Q emitter replaces the ordinary universal-distribution head
\[
 U(D,A)\longrightarrow
       \bigl(U(D,I(A,H))\longrightarrow U(D,H)\bigr)
\]
by
\[
 L_6^{\mathrm Q}(D,A,H)=
 U(D,I(A,H))\longrightarrow
       \bigl(U(D,A)\longrightarrow U(D,H)\bigr),
 \tag{Q}
\]
where each displayed implication is the numerical constructor $I$.
The other five ordinary heads, the dominated implicational head $\rho$,
the diagonal head $\Gamma$, and the guarded gate $G$ are unchanged from
the coefficient-four separated-diagonal calculus. The exchange formula
\[
 (P\longrightarrow(Q\longrightarrow R))\longrightarrow
       (Q\longrightarrow(P\longrightarrow R))
\]
gives proof translations in both directions. A retained 61-line
parametric $K,S,$ modus-ponens certificate is independently replayed in
this package. The availability of $K$ and $S$ from the unchanged restricted
implication basis remains its explicitly identified inherited dependency;
the exchange certificate does not assume the quantifier axiom it replaces.
The enclosing package additionally checks an expanded 18,897-row proof
directly from restricted $\rho_8$ and MP; its domination side conditions
are justified by syntactic subterm witnesses. This removes $K,S$ as
unverified premises of the relative-calculus argument.

\subsection{The numerical calculus, explicitly}

For completeness the remaining heads are fixed here, so the relative
statement does not depend on an unspecified axiom set. In the following
display each arrow abbreviates the numerical constructor $I$, not an
additional primitive logical rule. Put
$N(B)=I(B,1)$, $A\wedge B=N(I(A,N(B)))$, and
$\exists a\,B=N(U(a,N(B)))$. The six ordinary heads are
\[
\begin{aligned}
L_1(D,H)&=(\exists D\,U(D,H))\to H,\\
L_2(D,A,H)&=U(D,U(A,H))\to U(A,U(D,H)),\\
L_3(D,A,H)&=M(H,D)\to(E(H,A)\to M(A,D)),\\
L_4(D,A,H)&=E(H,D)\to(M(A,H)\to M(A,D)),\\
L_5(D,H)&=\exists D\,E(H,D),\\
L_6(D,A,H)&=U(D,A)\to(U(D,I(A,H))\to U(D,H)).
\end{aligned}
\]
The new calculus uses $L_6^{\mathrm Q}$ from (Q) in place of $L_6$.
Its additional heads are
\[
 \rho(D,A,H,d)=((H\to d)\to D)\to((D\to H)\to(A\to H)),
 \quad d\leq D,
 \qquad \Gamma(c)=c\to U(c,c).
\]
To define the gate without overloading the Q-variant name, let
\[
\begin{aligned}
 V(a,b,c)&=M(a,b)\to U(c,M(a,b)),\\
 P(a,b,c)&=c\to M(c,b),\\
 T(a,b,c)&=U\bigl(b,(M(a,b)\to P(a,b,c))\wedge(M(c,b)\to c)\bigr),\\
 R(a,b,c)&=U\bigl(c,T(a,b,c)\to N(P(a,b,c))\bigr),\\
 G(a,b,c)&=N\bigl(V(a,b,c)\to R(a,b,c)\bigr).
\end{aligned}
\]
The gate is supplied when $c\ne a$ and $c\ne b$; the emitted replacement
is zero otherwise. Modus ponens sends proofs of $A$ and $I(A,B)$ to a
proof of $B$. Theorem generalization sends a proof of $B$ to a proof of
$U(i,B)$ for any $i$. These, together with the displayed heads and their
specified side conditions, define the separate-diagonal numerical proof
predicate used in this report. A proof of zero is already available from
$\rho(0,0,0,0)=0$.

\subsection{The exact four-parameter command interface}

A command store is $(t,w,d,h)$: the visible theorem code, numerical stack,
arithmetic offset, and cached theorem code. Set $D=d+x$ and $H=h+y$.
The new axiom command has four parameters:
\[
 \mathsf A_{c,x,y,z}(t,w,d,h)
       =\bigl(0,[q_0,q_1,\ldots,q_9,t;w],z,0\bigr).
 \tag{A}
\]
The ten generated entries are as follows; the saved old top $t$ is the
\emph{eleventh} entry.
\begin{center}
\begin{tabular}{cl}
\toprule
Entry & Generated code\\
\midrule
$q_0$ & $G(H,D,c)$ if the numerical gate guard succeeds, otherwise $0$\\
$q_1$ & $\Gamma(c)$\\
$q_2$ & $L_6^{\mathrm Q}(D,c,H)$\\
$q_3$ & $L_5(D,H)$\\
$q_4$ & $L_4(D,c,H)$\\
$q_5$ & $L_3(D,c,H)$\\
$q_6$ & $L_2(D,c,H)$\\
$q_7$ & $L_1(D,H)$\\
$q_8$ & $\rho(D,c,H,d)$\\
$q_9$ & $U(d,h)$\\
\bottomrule
\end{tabular}
\end{center}
The dominated head has a valid side condition because $d\leq D=d+x$.
The gate guard is the inherited distinct-label condition on $c,H,D$.
All of $q_0,\ldots,q_8$ have residue zero modulo four. The last generated
entry $q_9$ has residue two and is at least two. Hence none of the generated
entries is the contradiction code one, for arbitrary incoming numerical
stores, independently of consistency.

The extra parameter $z$ is essential. After the gate has been pushed,
the emitter falls through the \emph{full} MP reader, which consumes a
further field of the current continuation. At that point $t=d=h=0$.
The ensuing MP attempt compares $I(0,z)$ with zero. It succeeds at $z=0$
and fails at $z>0$, but both cases return a visible zero, preserve the
new stack, clear the cache, and leave offset $z$. The proof compiler
chooses $z=0$. Soundness must cover every $z$, since the scheduler
enumerates every numerical candidate.

For $b=d+x$, the no-push MP command is exactly
\[
 \mathsf{MP}_x(t,w,d,h)=
 \begin{cases}
   (b,w,b,0),& I(t,b)=h,\\
   (0,w,b,0),& I(t,b)\ne h.
 \end{cases}
 \tag{MP}
\]
Successful MP does not insert the older dummy zero into $w$. Shift is
\[
 \mathsf S(t,w,d,h)=
 \begin{cases}
   \mathsf{halt},&h=1,\\
   \bigl(h,\operatorname{tl}(w),d,\operatorname{hd}(w)\bigr),&h\ne1.
 \end{cases}
 \tag{S}
\]
This is a change in recurring computation. For example, Shift halts at
$(t,h)=(1,1)$, whereas the earlier sum-one guard does not; a visible one
with cache zero is not directly tested. Changing these cases requires
changing the proof-command argument.

\subsection{The command identities that survive, and the necessary early halt}

Let $A_0=\mathsf A_{0,0,0,0}$ and define
\[
 \mathsf{Push}=A_0;\mathsf S^{10};\mathsf{MP}_0,
 \qquad
 \mathsf{Pop}=A_0;\mathsf S^{12},
 \qquad
 \mathsf{Discard}=\mathsf{MP}_0.
\]
Direct iteration of (A) and (S) gives
\[
 \begin{aligned}
 \mathsf A_{c,x,y,0};\mathsf S^{10}(t,w,d,h)
      &=(q_8,[t;w],0,q_9),\\
 \mathsf A_{c,x,y,0};\mathsf S^{11}(t,w,d,h)
      &=(q_9,w,0,t).
 \end{aligned}
\]
These shifts cannot halt: the tested generated entries are never one.
Thus Push is total and has effect
\[
       \mathsf{Push}(t,w,d,h)=(0,[t;w],0,0).
\]
Its last MP comparison fails because $q_9$ is a universal code and
$I(q_8,0)$ is an implication code. Pop has the conditional identity
\[
 \mathsf{Pop}(t,[a;w],d,h)=
 \begin{cases}
   \mathsf{halt},&t=1,\\
   (t,w,0,a),&t\ne1.
 \end{cases}
 \tag{Pop}
\]
The final Shift tests precisely the original visible theorem, now in the
cache. The old AX11 claim that Pop always preserves a visible one would
be false for this composition. The new proof compiler therefore permits
an early halt when it has already obtained contradiction.

The remaining elementary identities are
\[
 \mathsf{Discard}(t,w,0,h)=(0,w,0,0),
 \qquad
 \mathsf{MP}_i(0,w,0,0)=(0,w,i,0).
\]
For $0\leq j\leq8$, a generated head is selected by
\[
 \mathsf{Discard};\mathsf A_{A,D,H,0};
     \mathsf S^{j+2};\mathsf{Pop}^{9-j}.
\]
It returns $(q_j,w,0,0)$ from $(t,w,0,0)$. Every Pop preserves a
noncontradictory generated head. The restricted-head selector first
sets the actual offset $u$, then uses an increment $D-u$ with $u\leq D$;
this retains the numerical side condition rather than treating it as
an unrestricted axiom. Generalization is compiled by
\[
 \mathsf{Push};\mathsf{MP}_i;\mathsf S;
       \mathsf A_{i,0,0,0};\mathsf S^{11}.
\]
It sends $(a,w,0,0)$ to $(U(i,a),w,0,0)$, including $a=1$.

To combine premises $a$ and $I(a,b)$, the framed proof compiler
reserves a zero below them. The resulting local suffix is
\[
 \mathsf S;\mathsf S;\mathsf{MP}_b;\mathsf{Pop}
 \quad\hbox{on}\quad
 (0,[a,I(a,b),0;w],0,0).
\]
If $a=1$ it halts at the second Shift. If $a\ne1$ and $b=1$ it halts
inside Pop. Otherwise it returns $(b,w,0,0)$. These alternatives,
together with the separate all-zero calculation, cover the numerical
alias $[0;0]=I(0,0)=0$. If a compiled refutation ever returns a visible
one instead of halting earlier, the suffix
\[
                         \mathsf{Push};\mathsf S;\mathsf S
\]
exposes that one in the cache and halts. A lone final Shift would be
incorrect for the cached-only test.

The command checker proves the stated symbolic identities from total
Cantor projections, constructor residues, and explicit assumptions
$t\ne1$, $a\ne1$, or $b\ne1$ in the appropriate branches. It separately
checks each corresponding equality-to-one branch. It does not use a
consistency assumption to dismiss those cases.

\subsection{Scheduler encoding and source binding}

The four-field fresh code for AX is
\[
 \operatorname{enc}(\mathsf A_{c,x,y,z},e)
       =p\bigl(c+2,p(y,p(x,p(z,e)))\bigr).
\]
Positive continuation codes have the outer successor used by the
inherited scheduler. The independent checker verifies all four field
projections, including the collapsed-zero cases, and verifies that using
the obsolete three-field encoding would consume a different continuation.
For a universal halting suffix $W$, the fresh words
$\mathsf A_{k,0,0,0};W$ have a strictly increasing cofinal family of
candidate codes. The detailed arithmetic transfer argument must use
this actual scheduler and its carried store.

The finite source description is fully retained. Starting from the
pinned coefficient-four nodes, it changes exactly two named node
definitions: \texttt{node163} removes the successful-MP dummy push,
and \texttt{node170} implements the H guard. The axiom subtree is the
pinned Q grammar with its sole terminal \texttt{breakTo 1} replaced by
the existing stack-push macro \texttt{node49}. This macro then falls
through to the full MP reader. No other source-node definition changes.
Independent expansion of \texttt{node172} gives all 3,671 primitive sites.
Two independent control-flow resolvers agree on every instruction and
successor, and their result equals the delivered complete source graph.
A separate free-constructor calculation checks the entire AX stack,
the exchanged sixth head, the gate, and both gate-guard outcomes.

At the fixed-controller level, the new rows relative to the pinned
coefficient-four parent are
\begin{center}
\begin{tabular}{rrrcr}
\toprule
State & Read & Write & Move & Next\\
\midrule
38&0&0&R&11\\
38&1&1&R&122\\
164&1&1&R&114\\
214&0&0&R&312\\
214&1&1&R&20\\
312&0&0&R&188\\
\bottomrule
\end{tabular}
\end{center}
The Q reader graph and terminal AX push are generated separately from
the exact supplied grammar. All arithmetic-box transition rows remain
unchanged. The literal independent checker verifies those boxes from
rows rather than assuming their construction-time oracle.

\subsection{The quotient and the complete initial frame}

The raw machine has 134 fixed controls and 67 AX-reader controls.
Seven simultaneous history-sensitive folds give
\[
                         134+67-7=194.
\]
In native labels, the removed/retained pairs are
\begin{center}
\begin{tabular}{rr@{\qquad}rr}
\toprule
Removed&Retained&Removed&Retained\\
\midrule
2004&2052&2021&2030\\
2025&2035&2027&2046\\
2039&2044&2055&2056\\
2058&2017&&\\
\bottomrule
\end{tabular}
\end{center}
For each pair, every incoming transition to the first member writes
zero and moves right; every incoming transition to the second writes
one and moves right. The preceding tape cell therefore identifies the
old control. Four complete two-cell cases per pair verify the quotient
on the final simultaneously folded table. Each case recovers the
aliased raw successor, the raw head displacement, the raw written
symbol, the unchanged history cell, and arbitrary untouched exterior
tape. All 28 cases pass, with durations one, three, or five steps.
This verifies compatibility of the seven folds together.

A breadth-first renumbering from native state 74 visits all 194 working
controls and maps halt 254 to 194. The dense/native conjugacy is checked
independently. Starting from an entirely blank tape, the machine reaches
the complete clean source frame after exactly 57,922 physical steps:
native state 84, head coordinate $-2$, the address of source site 25,
header base 48, logical registers $(2,0,\ldots,0)$, and 26 spare unary
blocks of value one. Both startup implementations check the entire
finite-support tape, not just the visible register values; all exterior
cells remain zero. Their tape representations are independent: a set
of one-cells in the construction checker and a byte array in the literal
checker.

\subsection{Observed verification and exact deliverable identity}

\begin{center}
\begin{tabular}{lr}
\toprule
Quantity & Observed value\\
\midrule
Working states / transition rows&194 / 388\\
Primitive source sites&3,671\\
Universal source outcomes&4,778\\
Construction symbolic operations&1,547,747\\
Independent symbolic operations&1,130,254\\
Independent maximum case length&403\\
Minimum positive duration constant&3\\
Generic list lemmas rederived from literal rows&4\\
Complete blank initialization steps&57,922\\
\bottomrule
\end{tabular}
\end{center}

The independent physical contracts quantify over nine arbitrary natural
register values, both decrement outcomes, an arbitrary nonempty finite
list of arbitrary natural spare values, and an affine header gap. The
independent engine imports neither the construction compiler nor its
arithmetic-box oracle. A fresh aggregate offline run rebuilt every table
and source artifact, checked source binding, checked the symbolic command
identities, replayed the premise-exchange certificate, and reran this
independent literal-tape engine. Every stage passed; the full observed
receipt is \texttt{verification\_full.json}.

The dense table's SHA-256 is
\begin{quote}\small\ttfamily
19e215df20429fa6a35f0a99ba0dcfcb40893a6e4b9a4ae6e450b0225f110e50
\end{quote}
The native table's SHA-256 is
\begin{quote}\small\ttfamily
424bd5175bbca161e76af485f1eda03f2f169d90afe80563035b60289361bf07
\end{quote}
The tables, source graphs, fold witnesses, command certificates, preserved
licences, input hashes, and reconstruction code are all included. Running
\texttt{python3 -B verify.py} reruns both physical engines in an isolated
temporary directory. Its optional \texttt{--quick} mode skips only the
independent physical replay and explicitly records that omission.

\subsection{New search dimension: permuting emitted theorem heads}

The seven separately emitted heads---the six ordinary logical heads and
the diagonal head---can be permuted while preserving the numerical
calculus. The command selectors must follow the new stack positions.
This supplies a search dimension beyond reassociation of a fixed macro
word. The retained search combines binary reassociation with pairwise
exchanges of these seven blocks, retaining the initial generalization
and dominated-axiom prefix and the final guarded gate. Near the target
frontier it checks exact compatibility of simultaneous two-cell history
folds; it does not certify a score from matching alone.

The first 240-second run, using the supplied AX11 controller, considered
3,631,115 proposals and accepted 650,390 search moves. The second
90-second run used the new H-controller rows, making their different
provider transitions available to decoder folds; it considered
1,283,873 proposals and accepted 221,857 moves. Neither run obtained an
effective reader count below 60, hence neither obtained a machine below
194 states with its 134 fixed controls. Exact inputs, seeds, source,
logs, and receipts are included. These are bounded heuristic searches,
not a general state-count lower bound, and no unverified numerical score
has been promoted to a smaller machine.


\section{Independent verification and its proof boundary}
\label{sec:verification}

The full candidate verifier rebuilds the table from pinned inputs in a
fresh temporary directory, verifies exact equality of the delivered source
artifacts, resolves control flow by two independent algorithms, checks the
free-term emitter and command identities, and replays the literal-tape
checker. The last checker does not import the construction compiler or its
arithmetic-box oracle. It rederives four generic finite-list lemmas from
literal transition rows and applies them to every source outcome.

The quantification is substantial: all nine register values are arbitrary
naturals, both zero and positive decrement outcomes are checked, the spare
list is an arbitrary nonempty finite list of arbitrary natural values, and
the header gap is $48+n$ for arbitrary $n\geq0$, less the source-address
length. This is stronger than testing a finite collection of small register
values. Its remaining trust consists of the mathematical correctness of
the symbolic rules and their executable implementation; it is not a newly
compiled Lean or Rocq theorem.

The blank initialization is checked separately, by a byte-array simulator
independent of the construction's set-of-ones simulator. The entire finite
support is compared, including left marker, exact address, gap, all nine
unary registers, the 26 spare blocks, final head, and zero exterior. The
computation has not halted before this frame. A finite simulation is used
to prove this specific initialization fact, not to infer nonhalting.

For the old primary PA194, the same literal-row engine is instantiated with
the exact pinned dense and native tables, its 3,673 source sites, 4,780
outcomes, and 58,154-step startup. The separately saved fresh replay binds
the relative argument to the actual old table rather than to a similarly
named machine. The comparison's historical status labels are not used as
substitutes for this replay.

Three levels must stay separate. A concrete symbolic execution certificate
establishes its specified finite or universally parameterized machine
contract. A written metamathematical argument explains why the coded
contract can be proved using the restricted induction of iPRA. A completed
proof-assistant development would also formally check that explanation and
its arithmetic encoding. The deliverable provides the first level and
written arguments at the second, with the precise remaining calibration
obligation stated below. It does not claim the third.

% A standalone mathematical section.  This is not a compiled proof-assistant theorem.
% Requires amsmath, amssymb and amsthm in the containing document.
\section{What an iPRA calibration must establish}
\label{sec:ipra-audit}

Let $\operatorname{Prf}_{\rm PA}(p)$ be the primitive-recursive predicate
that $p$ is a PA proof of $0=1$.  For a fixed literal Turing table $M$,
let $H_M(t)$ be the primitive-recursive predicate that the absorbing
totalization of its blank run is in the halt state at time $t$.  We use
\[
  \operatorname{Con}({\rm PA})\equiv
       \forall p\,\neg\operatorname{Prf}_{\rm PA}(p),\qquad
  \operatorname{NH}(M)\equiv\forall t\,\neg H_M(t).
\]
All the predicates in this paragraph may be represented by their
characteristic-function equations in iPRA.  This is a claim about the
\emph{metatheory in which equivalence is proved}; the theory whose proofs
are searched remains classical PA.

\subsection{A direct constructive criterion}

\begin{proposition}[Witness-transformer criterion]
Suppose $f$ and $g$ are primitive-recursive functions and iPRA proves
\begin{align}
 \operatorname{Prf}_{\rm PA}(p)&\longrightarrow H_M(f(p)),
       \label{eq:pr-to-halt}\\
 H_M(t)&\longrightarrow\operatorname{Prf}_{\rm PA}(g(t)).
       \label{eq:halt-to-pr}
\end{align}
Then iPRA proves
\[
 \operatorname{NH}(M)\ \longleftrightarrow\ \operatorname{Con}({\rm PA}).
\]
\end{proposition}
\begin{proof}
Under nonhalting, a PA proof $p$ would yield $H_M(f(p))$ by
\eqref{eq:pr-to-halt}, a contradiction.  Conversely, under consistency,
a halting time $t$ would yield a PA proof $g(t)$ by
\eqref{eq:halt-to-pr}, also a contradiction.  These are intuitionistic
implication introductions; neither Markov's principle, a reflection
principle, nor excluded middle for an unbounded formula is used.
\end{proof}

The hypothesis is not just that $f$ and $g$ happen to be primitive
recursive or that PA proves their correctness.  The displayed correctness
statements themselves must be available in iPRA.  A proposed compiler
with correctness proved using unrestricted induction does not yet meet
this criterion.

\subsection{A constructive finite-segment transfer theorem}

The physical part of a state reduction has a useful sufficient condition
which is considerably weaker than reproving the PA calibration.
Let $s$ and $m$ be primitive-recursive, deterministic step functions on
coded source and machine configurations.  Halting configurations are
absorbing.  Let $h_s$ and $h_m$ be their primitive-recursive halt tests.
Assume a primitive-recursive frame map $F$, a primitive-recursive invariant
$J$, and a primitive-recursive positive length $\ell$ satisfy the following
formulas, with proofs in iPRA:
\begin{align*}
 &J(c_0),\qquad J(c)\longrightarrow J(s(c)),\\
 &J(c)\longrightarrow 1\leq\ell(c),\\
 &J(c)\longrightarrow m^{\ell(c)}(F(c))=F(s(c)),\\
 &J(c)\longrightarrow\bigl(h_m(F(c))\leftrightarrow h_s(c)\bigr).
\end{align*}
Here $s^n$ and $m^n$ denote primitive-recursive iteration.  A halt source
configuration may be given length one; absorption then supplies its
segment equation.  Assume also a fixed blank initialization of length
$b$ ending at $F(c_0)$, with a verified nonhalting prefix.  A source which
is already halted at $c_0$ can instead be handled as a separate fixed case.

\begin{proposition}[Positive-segment transfer in iPRA]
Under these hypotheses iPRA proves that the source computation and the
literal blank-start machine halt under exactly the same circumstances.
Both directions have primitive-recursive witness transformations.
\end{proposition}
\begin{proof}
Define simultaneously by primitive recursion
\[
 c_n=s^n(c_0),\quad
 L_0=0,\quad L_{n+1}=L_n+\ell(c_n).
\]
Quantifier-free induction, with the primitive-recursive predicates and
finite-word operations replaced by their function symbols, gives
\[
 J(c_n),\qquad L_n\geq n,\qquad
 m^{b+L_n}(\text{blank})=F(c_n).
\]
In particular a source halt at $n$ gives a machine halt at $b+L_n$.

For the reverse direction suppose the machine has halted at time $t$.
Perform the bounded search for the first $n\leq t+1$ with $h_s(c_n)$;
this is primitive recursive.  If the search succeeds it supplies the
required source witness.  Otherwise its decidable search certificate
gives $\neg h_s(c_{t+1})$.  The displayed simulation gives a nonhalting
machine frame at time $b+L_{t+1}>t$.  Absorption of machine halting,
proved by quantifier-free induction on the additional number of steps,
contradicts that frame.  Since the search has a decidable result, this
contradiction proves that its success branch holds and yields its
bounded witness intuitionistically.  The initialization-prefix check
also excludes a halt preceding $b$.

All inductions used here concern equations or decidable predicates of
explicitly defined primitive-recursive functions.  The proof does not
induct on ``there exists a next frame'' and does not use an unbounded
least-time search or a standard-model soundness argument.
\end{proof}

If two literal machines satisfy such contracts for the same source, the
proposition gives an iPRA proof of their halting equivalence.  It then
transfers any iPRA consistency calibration already proved for either
machine.  It does not manufacture a missing PA-to-source proof compiler.
For a candidate-specific application, the actual $J,F,\ell$, their
primitive-recursive definitions, every successor-frame contract and the
blank initialization must be supplied.  Finite numerical tests of some
register values do not establish these hypotheses.

\subsection{The alternative I\texorpdfstring{$\Sigma_1$}{Sigma1} route}

One can instead prove, in $\mathrm I\Sigma_1$,
\[
 \exists p\,\operatorname{Prf}_{\rm PA}(p)
       \ \longleftrightarrow\ \exists t\,H_M(t).
\]
The two witness implications can be put in $\Pi_2$ form.  The
$\Pi_2$ conservation of $\mathrm I\Sigma_1$ over classical PRA then
gives the corresponding classical PRA implications.  Applying the
G\"odel--Gentzen interpretation and using decidability of the open
matrices gives the desired equivalence of the two negative sentences
in iPRA.  No elimination of the double negation of a $\Sigma_1$ sentence
is needed.  ProveIt's \texttt{isigma1-transfer} and
\texttt{isigma1-machine} directories formalize this route, conditional
on the machine-specific calibration premise.

The relevant distinction is between a theorem about the ordinary natural
numbers and a proof in the weak theory.  An external proof-assistant
theorem about all standard finite runs can support the finite
combinatorics, but it is not by itself the required internal arithmetic
proof.  Conversely, the absence of a completed Lean internalization does
not by itself disprove the possibility of a written iPRA proof: an
explicit primitive-recursive compiler with verified decidable invariants
would suffice.

\section{A relative equivalence in intuitionistic primitive recursive arithmetic}
\label{sec:relative-ipra}

This section proves a relative statement about the delivered machine and
the primary PA194 table.  It does not assume that the primary table has
already been calibrated to PA in iPRA, and it does not prove that missing
calibration.  The proof here is an ordinary mathematical proof with finite
symbolic certificates; it is not a newly compiled Lean or Rocq theorem.

Write $M_0$ for the primary table with SHA-256
\path{4bfea62b17aad7d8c8c5e682fe8942e0a86b93e4351f8e675e32705bdaae007f}
and $M_1$ for PA194AX11QH, whose table has SHA-256
\path{19e215df20429fa6a35f0a99ba0dcfcb40893a6e4b9a4ae6e450b0225f110e50}.
In this section a halting sentence means an existential assertion about
the primitive-recursive iteration of the literal blank-start table,
with the halt state made absorbing.  Thus no external truth predicate
or interpretation of arbitrary PA formulas occurs in the assertion.

\begin{theorem}[Relative machine equivalence]
\label{thm:relative-ipra}
For the two specified tables,
\[
 \iPRA\vdash \Halt(M_0)\leftrightarrow\Halt(M_1),
 \qquad
 \iPRA\vdash \NH(M_0)\leftrightarrow\NH(M_1).
\]
Consequently any iPRA calibration of either table to $\Con(\PA)$
would transfer to the other.  This consequence does not discharge the
calibration premise.
\end{theorem}

The details below explain why the inductions needed for this relative
statement are available in iPRA, instead of merely asserting that all
the algorithms are primitive recursive.

\subsection{The two finite numerical proof predicates}

For clarity, all numerical constructors used here are fixed explicitly.
Let $T(0)=0$, $T(s+1)=T(s)+s+1$, and $p(a,b)=T(a+b)+b$.
Write
\begin{align*}
 I(a,b)&=4p(a,b),& M(a,b)&=4p(a,b)+1,\\
 U(a,b)&=4p(b,a)+2,& E(a,b)&=4p(a,b)+3,\\
 N(a)&=I(a,1),& C(a,b)&=N(I(a,N(b))),\\
 X(a,b)&=N(U(a,N(b))).
\end{align*}
The six heads, in the parameter order used by the emitter, are
\begin{align*}
 L_1(D,H)&=I(X(D,U(D,H)),H),\\
 L_2(D,A,H)&=I(U(D,U(A,H)),U(A,U(D,H))),\\
 L_3(D,A,H)&=I(M(H,D),I(E(H,A),M(A,D))),\\
 L_4(D,A,H)&=I(E(H,D),I(M(A,H),M(A,D))),\\
 L_5(D,H)&=X(D,E(H,D)),\\
 L_6(D,A,H)&=I(U(D,A),I(U(D,I(A,H)),U(D,H))),\\
 L_{6Q}(D,A,H)&=I(U(D,I(A,H)),I(U(D,A),U(D,H))).
\end{align*}
The dominated head and diagonal are
\[
 \rho_8(a,b,c,u)=I(I(I(c,u),a),I(I(a,c),I(b,c))),
 \qquad \Gamma(c)=I(c,U(c,c)).
\]
To define the guarded gate, put
\begin{align*}
 V(a,b,c)&=I(M(a,b),U(c,M(a,b))),\\
 P(a,b,c)&=I(c,M(c,b)),\\
 R(a,b,c)&=U\bigl(b,C(I(M(a,b),P(a,b,c)),I(M(c,b),c))\bigr),\\
 R'(a,b,c)&=U(c,I(R(a,b,c),N(P(a,b,c)))),\\
 G(a,b,c)&=N(I(V(a,b,c),R'(a,b,c))),\\
 \widehat G(a,b,c)&=
   \begin{cases}0,&c=a\text{ or }c=b,\\G(a,b,c),&\text{otherwise.}\end{cases}
\end{align*}
This is the separate-diagonal gate called \texttt{oldBundle} in the
pinned source; it is not the different equality-vacuity bundle also
defined there.  The generated entries are, exactly,
\begin{align*}
 (q_0,q_1,q_2,q_3,q_4)
   &=\bigl(\widehat G(H,D,c),\Gamma(c),L_{6Q}(D,c,H),L_5(D,H),L_4(D,c,H)\bigr),\\
 (q_5,q_6,q_7,q_8,q_9)
   &=\bigl(L_3(D,c,H),L_2(D,c,H),L_1(D,H),\rho_8(D,c,H,d),U(d,h)\bigr),
\end{align*}
with $L_6$ in place of $L_{6Q}$ for the primary machine.

Let $\mathcal C$ be the separate-diagonal numerical calculus of the
primary machine.  It has the six displayed heads, the dominated
implication head $\rho_8$, the guarded gate, the diagonal clause, modus
ponens and theorem generalization.  Let $\mathcal C_Q$ change only
\[
 P\to(Q\to R)\quad\hbox{to}\quad Q\to(P\to R)
\]
in the sixth head.  Here $P=U(D,A)$, $Q=U(D,I(A,H))$, and $R=U(D,H)$.
More precisely, an axiom proof is an instance of $L_1,\ldots,L_6$,
of $\rho_8(a,b,c,u)$ with $u\leq a$, of $G(a,b,c)$ with
$c\ne a,b$, or of $\Gamma(c)$.  The two rules are
$(a,I(a,b))\mapsto b$ and $a\mapsto U(i,a)$ for arbitrary $i$.
This fixes the finite proof checker without appealing to a semantic
interpretation of these numerical expressions.

Proofs can be coded as finite trees with their conclusion included at
each node.  Use a tagged pairing code which is larger than the sum of
its child codes and its data fields.  The assertion that such a tree
is valid, its conclusion, its height, and its number of nodes are all
primitive recursive.  For a valid proof code $D$, every child code,
displayed conclusion, and numerical parameter is at most $D$, and the
height and number of nodes are at most $D$.  These properties are
proved by bounded tree traversal and quantifier-free induction.
Every leaf has height one and an internal node has height one plus
the maximum child height; there are no valid height-zero proofs.

The exchange certificate has 29,932 free implication terms and 18,897
proof rows: 8,986 instances of the \emph{restricted} $\rho_8$ axiom
and 9,911 modus-ponens rows.  The independent check verifies
\[
 (a\to(b\to c))\to(b\to(a\to c)).                 \tag{E}
\]
Each domination condition $u\leq a$ in a $\rho_8$ row is witnessed
by the fact that its term $u$ is a syntactic subterm of its parameter
$a$.  Since $I(x,y)=4p(x,y)$ is at least both $x$ and $y$, each such
inequality has an iPRA proof.  Each MP row is an identity in the free
implication algebra.  Instantiating all 18,897 rows at arbitrary natural
$a,b,c$ therefore gives a fixed iPRA-verifiable proof-producing template.
The certificate does not assume $K$, $S$, unrestricted $\rho_8$, or
quantifier distribution as additional axioms.

Replace a sixth-head leaf by the opposite sixth-head leaf and (E),
followed by MP; keep every other leaf and both proof constructors.
Doing this in either direction is primitive recursive.  Its validity
is verified by bounded induction through the coded input tree: all
intermediate certificates are finite lists, every local validity test
is decidable, and the fixed exchange template contributes no unbounded
proof search.  Thus iPRA proves
\[
 \exists e\,\operatorname{Prf}_{\mathcal C}(e,1)
 \quad\longleftrightarrow\quad
 \exists e\,\operatorname{Prf}_{\mathcal C_Q}(e,1).       \tag{1}
\]
The numerical alias $I(0,0)=0$ is harmless: evaluation of a symbolic
identity preserves it, and the MP predicate tests numerical equality.

\subsection{The exact command interfaces and finite words}

Put $[a;w]=p(a,w)$, with total projections $\operatorname{hd}$ and
$\operatorname{tl}$.  For the new machine the commands, on the store
$(t,w,d,h)$, are
\begin{align*}
 S(t,w,d,h)&=
 \begin{cases}\mathsf{halt},&h=1,\\
 (h,\operatorname{tl}w,d,\operatorname{hd}w),&h\ne1,
 \end{cases}\\
 \operatorname{MP}_x(t,w,d,h)&=
 \begin{cases}(b,w,b,0),&I(t,b)=h,\\(0,w,b,0),&I(t,b)\ne h,
 \end{cases}\qquad b=d+x,\\
 A_{c,x,y,z}(t,w,d,h)&=
       (0,[q_0,q_1,\ldots,q_9,t;w],z,0).
\end{align*}
Here $D=d+x$, $H=h+y$, and the generated entries are the gate,
diagonal, exchanged sixth head, the other five heads, the dominated
head, and $U(d,h)$, in that order.  They are exactly the entries in
the free-constructor emitter certificate.  In particular $q_0,\ldots,q_8$
are zero modulo four and $q_9$ is two modulo four, so none is one.

Let $A_0=A_{0,0,0,0}$.  The following finite words have the indicated
effects for \emph{arbitrary} numerical stacks, including $[0;0]=0$:
\[
 \operatorname{Push}=A_0S^{10}\operatorname{MP}_0,
 \qquad\operatorname{Pop}=A_0S^{12},
 \qquad\operatorname{Discard}=\operatorname{MP}_0.
\]
Push sends $(t,w,d,h)$ to $(0,[t;w],0,0)$.  Discard sends
$(t,w,0,h)$ to $(0,w,0,0)$.  Pop, on $(t,[a;w],d,h)$, halts if
$t=1$ and otherwise returns $(t,w,0,a)$.  All these statements follow
by a fixed number of substitutions in the displayed command equations.
They are quantified equations with decidable case splits and require
no induction over formula truth.

For $0\leq j\leq8$, the selector
\[
 \operatorname{Discard};A_{A,D,H,0};S^{j+2};
                  \operatorname{Pop}^{9-j}                 \tag{2}
\]
returns $(q_j,w,0,0)$ from any $(t,w,0,0)$.  For the dominated
instance $\rho_8(D,A,H,u)$, with $u\leq D$, use
\[
 \operatorname{Discard};\operatorname{MP}_u;
       A_{A,D-u,H,0};S^{10};\operatorname{Pop}.             \tag{3}
\]
The old cache is zero after Discard; consequently the intermediate
MP sets the offset to $u$, including at $u=0$.  Generalization is
\[
 \operatorname{Gen}_i=\operatorname{Push};\operatorname{MP}_i;S;
                A_{i,0,0,0};S^{11},                       \tag{4}
\]
which returns $(U(i,a),w,0,0)$ on $(a,w,0,0)$, even when $a=1$.

On $(0,[a,I(a,b),0;w],0,0)$, the suffix
$S;S;\operatorname{MP}_b;\operatorname{Pop}$ either halts
(when $a=1$ or $b=1$) or returns $(b,w,0,0)$.  Therefore, at an
MP proof node with major proof $e_1$ and minor proof $e_0$, compile
\[
 W_e=\operatorname{Discard};\operatorname{Push};W_{e_1};
       \operatorname{Push};W_{e_0};\operatorname{Push};
                 S;S;\operatorname{MP}_b;\operatorname{Pop}.  \tag{5}
\]
The initial Push reserves the zero needed by the suffix.  Axiom leaves
use (2)--(3), and a generalization node appends (4).  It follows that
\[
 \operatorname{Eval}(W_e,(t,w,0,0))=\mathsf{halt}
 \quad\text{or}\quad
 \operatorname{Eval}(W_e,(t,w,0,0))=(\operatorname{concl}(e),w,0,0).
                                                               \tag{6}
\]
Thus, when $\operatorname{concl}(e)=1$, the word
\[
             V_e=A_0;W_e;\operatorname{Push};S;S            \tag{7}
\]
halts from every carried store.  A lone final Shift would not suffice
for the new cache-only guard.

For completeness, the primary machine uses the same arithmetic heads
before the exchange, but its AX leaves $q_0$ visible and pushes
$q_1,\ldots,q_9,t$; its successful MP additionally pushes a zero;
its Shift tests $t+h=1$.  Its corresponding words are
\begin{align*}
 \operatorname{Push}_0&=A_0S^9\operatorname{MP}_0,&
 \operatorname{Pop}_0&=A_0S^{11},\\
 \operatorname{Discard}_0&=\operatorname{MP}_1;
       A_{1,0,0};S^{10};\operatorname{MP}_0,&
 \operatorname{Offset}_i&=
      \begin{cases}\epsilon,&i=0,\\\operatorname{MP}_i,&i>0.
      \end{cases}
\end{align*}
Here the primary Discard contract is asserted on inputs with
$d=h=0$, exactly as supplied by its compiler; the new Discard's
arbitrary-cache contract is not attributed to the primary word.
Its Pop is total even when $t=1$.  The $q_j$ selector, $1\leq j\leq8$,
uses $\operatorname{Discard}_0;A_{A,D,H};S^{j+1};
\operatorname{Pop}_0^{9-j}$, with (3)'s analogous offset treatment
for a nonzero lower bound.  Its gate selector is
$\operatorname{Discard}_0;A_{c,b,a};\operatorname{Pop}_0^{10}$.
Generalization uses Push, Offset, Shift, $A_{i,0,0}$, and ten Shifts.
At an MP node there is no initial reserved-zero Push: the successful
MP supplies that zero.  The suffix is
$S;S;\operatorname{MP}_b;A_{b,0,0};S^{11}$.
The primary refutation word is $A_0;W_e;S$.
Direct substitutions give the same invariant (6) for its compiler.
These exact identities are also displayed in the repository's
\texttt{PA194Words.lean} and \texttt{PA194DerivationCompiler.lean};
the argument here uses their equations, not an unproved transfer of
an external Lean induction into iPRA.

\subsection{Why the compiler induction is weak enough}

The phrase ``induct on proofs, uniformly in the old stack'' can hide
an induction with unbounded quantifiers.  Here it can be replaced by
an explicitly bounded induction.  Fix a decorated proof-code bound $D$
and an initial bound $B_0\geq\max(D,t,w,1)$.  Define
\[
             B_{j+1}=(4B_j+4)^8.                          \tag{8}
\]
This bounds \emph{inputs of recursive subproof calls}, not every
intermediate value of an AX execution.  Before such a call in (5),
at most two stack pushes have been made, their pushed conclusions
are at most $D$, and each push satisfies
\[
        p(a,w)\leq(a+w+1)^2.
\]
Thus (8) is more than sufficient.  The primary compiler needs no
larger bound.  A successful subproof restores its protected old
stack, so its internal AX values do not enter this bound.

Put $k^*=\min(k,D)$.  Let $P(k)$ say that (6) holds for every valid proof
code at most $D$ of height at most $k^*$, and for all
$t,w\leq B_{D-k^*}$.
All these quantifiers are bounded, and evaluation of a finite word is
primitive recursive.  Thus $P(k)$ has a primitive-recursive characteristic
function and is usable as a quantifier-free induction predicate in
iPRA.  The base case is empty.  At a step from $k$ to $k+1$ with $k<D$, the
recursive calls have height at most $k$ and inputs at most $B_{D-k}$,
by (8).  The fixed leaf and suffix identities prove the remaining
cases.  If a proof has smaller height already, monotonicity of the
bounds applies.  For $k\geq D$ the predicate is unchanged, so those
induction steps are immediate.  At $k=D$ the result is (6) for the originally arbitrary
$t,w$.  This supplies the required uniformity without induction on an
unbounded $\forall w$ formula.

The compiler has at most 256 commands per input tree node, so its
word length is at most $256D$.  This bound is deliberately loose:
the largest leaf selector has fewer than 140 commands, and the
nonrecursive overhead at an MP or generalization node is smaller.
Repeated primitive-recursive concatenation therefore defines the
compiler and its complete evaluator.  No proof search is involved.

\subsection{Soundness by explicit proof annotations}

An annotation consists of a valid numerical proof of $t$, a valid
proof of $h$, and a \emph{tagged finite list} of valid proof codes for
the stack.  Fold the conclusions of that list with $p$ and tail zero
to obtain the numerical stack $w$.  The annotation need not be the
canonical Cantor expansion: a trailing theorem of zero is permitted.
For an empty annotation list, use the fixed proof of zero, which is
the instance $\rho_8(0,0,0,0)=0$.

This choice makes all zero aliases exact.  Pushing a zero proof can
extend the tagged annotation list even when $p(0,0)=0$ leaves the
numerical stack unchanged.  For a nonempty list the Cantor projection
identities give the corresponding head and tail; for an empty list
they give zero and zero.  No fictitious nonzero list length is read
off the number zero.

AX attaches its finite axiom proofs, a generalization of the old
cache proof, and the saved top proof.  A failed guard attaches the
zero proof.  Successful MP combines its two premise proofs; failure
attaches zero.  A primary successful MP also attaches the dummy-zero
proof to the stack.  Shift moves proofs between these positions.
The offset has no proof obligation.  The new halt test reads a proof
of one in the cache.  For the primary test $t+h=1$, a decidable
nonnegative-integer case split finds a proof of one in the top or
cache.

These updates are primitive recursive, including the updating of
the list of proof codes.  Their invariant is the decidable statement
``all listed proof objects pass the proof checker and their folded
conclusions equal the displayed store''.  Quantifier-free induction
on the scheduler clock therefore constructs, from any finite halting
scheduler run, an actual numerical proof of one.  This is a syntactic
proof transformer; it assumes no consistency or soundness of either
numerical calculus.

\subsection{An explicit bound for scheduler fairness}

For a finite word use $E(\epsilon)=0$ and
$E(c::v)=1+b_c(E(v))$.  The new command constructors are
\[
 b_S(z)=p(0,z),\quad
 b_{\operatorname{MP}_x}(z)=p(1,p(x,z)),\quad
 b_{A_{c,x,y,u}}(z)=p(c+2,p(y,p(x,p(u,z)))).                \tag{9}
\]
The primary AX has the same expression with its extra $p(u,z)$
replaced by $z$.  A fresh candidate $n$ is decoded directly; a
positive continuation $l$ is decoded from $l-1$.  Every extracted
continuation is at most its input seed.  Hence a positive continuation
strictly decreases, and candidate $n$ takes at most $n+1$ commands
unless it halts.

Given the all-store refutation word $V_e$, choose the single candidate
\[
 n=p(2,p(0,p(0,p(0,E(V_e)))))\geq3                        \tag{10}
\]
for the new machine; omit the innermost $p(0,\cdot)$ for the primary.
It first executes $A_0$ and then $V_e$.  The prefix may alter the
carried store, since (7) works from every store.  The fresh counter
starts at two and increases by one.  Therefore, by fewer than
\[
                  (n+2)^2                               \tag{11}
\]
command steps, the scheduler has either halted earlier or executed
this halting candidate.  Indeed, the commands in candidates
$2,\ldots,n$ number at most $\sum_{j=2}^n(j+1)<(n+2)^2$.
The comparison and the decreasing-continuation argument are bounded
induction on natural numbers and decidable run predicates.  Thus
(10)--(11) provide a primitive-recursive halting witness, rather
than an unbounded appeal to eventual enumeration.

Combining annotations with the two compilers proves in iPRA
\[
 \Halt(\text{primary scheduler})\leftrightarrow
        \exists e\,\operatorname{Prf}_{\mathcal C}(e,1),
 \qquad
 \Halt(\text{new scheduler})\leftrightarrow
        \exists e\,\operatorname{Prf}_{\mathcal C_Q}(e,1).   \tag{12}
\]

\subsection{Primitive source routines: explicit recurrences}

The source-to-command step is also witnessed by primitive-recursive
traces.  Here are its precise cutpoints and workspace invariant.
At the root the nine registers, in their source order, are
\[
       (R_0,\ldots,R_8)=(n,t,w,0,d,h,0,0,l).
\]
At an AX entry only $R_3=a$ is different.  Write $X=R_6$ and
$Y=R_7$ for the two work registers.  All ordinary arithmetic helper
contracts enter and leave with $X=Y=0$ and preserve every register
not mentioned in their effect.  There is one deliberate exception:
the two halves of the gate have the intermediate workspace
$(R_3,X,Y)=(0,a,a)$, specified below.  Source cutpoints are
\[
\begin{array}{c|rrrrr}
 &\text{root}&\text{AX}&\text{MP}&\text{return jump}&\text{Shift}\\ \hline
 M_0&25&50&3554&3660&3661\\
 M_1&25&50&3572&3659&3660
\end{array}
\]
The return jump has successor 25.  The shared prefix 0--24 contains
the explicit halt at 10 and stack unpairing at 11--24.  Fetch and
dispatch occupy 25--49.  These are source ordinals, not physical
Turing-machine state numbers.

The finite checker \texttt{build\_macro\_plan.py} records all 190 AX
token intervals.  It checks their exact resolved instructions and
all successor labels, that each interval has only its first site
as an external entry, and that its only external exit is the next
token's first site.  Of these intervals, 57 are the literal
instruction incrementing $R_1$; the other 133 match one of nine
pinned primary helper graphs under translation of their internal
labels and exit.  Their contracts are as follows; fields omitted
from a displayed assignment remain fixed.
\[
\begin{array}{c|l}
 \text{source node}&\text{effect from a clean workspace}\\ \hline
 49&(t,w)\mapsto(0,p(t,w))\\
 55&(t,w)\mapsto(d,p(t,w))\\
 58&(t,w)\mapsto(h,p(t,w))\\
 70&(t,w)\mapsto(a,p(t,w))\\
 65&(t,p(v,w))\mapsto(I(v,t),w)\\
 72&(h,l)\mapsto(h+\operatorname{hd}l,\operatorname{tl}l)\\
 75&(t,d,l)\mapsto(I(t,d+\operatorname{hd}l),
                      d+\operatorname{hd}l,\operatorname{tl}l)\\
 142&(a,X,Y)=(a,0,0)\mapsto(0,a,a)
\end{array}
\]
The remaining node 150 has the contract
\[
 (t,d,h,X,Y)=(t,d,h,a,a)\longmapsto(g,0,0,0,0),
 \qquad
 g=\begin{cases}0,&h=a\text{ or }d=a,\\t,&\text{otherwise,}\end{cases}
\]
with $R_3=0$ throughout.  Nodes 55, 58 and 70 are
\emph{push followed by copy}, not mere additions to $t$.
Node 75 reads $d$, pushes the incoming $t$, copies the updated $d$,
and constructs $I(t,d)$; its temporary push is consumed and $w$
is preserved.  The contracts hold for arbitrary numerical $w$,
including zero aliases.  At every use of node 65 the symbolic
emitter certificate supplies a nonempty \emph{tagged} stack prefix;
the equation $p(v,w)$, not a purported numerical list length,
justifies taking its head.  Node 75 supplies its own pushed item.
Their detailed row-cycle derivations are
the pinned \texttt{AXMacroSound}, \texttt{Readers},
\texttt{AXGateCompare} and \texttt{PA194AXGate} modules, with the
following explicit primitive-recursive clocks and induction
interpretation.  Equality of every resolved helper row, including
its exit, is what permits transplantation of those derivations;
an abstract macro name alone would not do so.

The independent \texttt{check\_outside\_ax\_rows.py} checks all
140 retained rows outside AX and every successor under the explicit
relocation.  It removes exactly the old MP dummy-push interval
$[3634,3653)$, redirects success to the retained copy-$d$ block at
3652, and checks the two new Shift blocks against the old clear-$t$
and copy-$h$ blocks.  These 140 rows, nine Shift-prefix rows and
the AX intervals cover all 3,671 new source sites.  Thus the two
finite binding checks together cover the complete source, including
the command dispatcher, not only the constructor expression.

The arithmetic routines have only
draining/copying loops, Cantor pairing and unpairing, multiplication
by four, and destructive comparisons.  These are not unspecified
arithmetic oracles.  For example, the actual pairing loop has length
\[
 P(0)=1,\qquad P(s+1)=P(s)+7s+5.
\]
The actual pairing core on inputs $t,x$ takes $4t+1+P(t+x)$ steps,
and pushing $t$ onto $w$ takes
\[
        3t+7w+4+P(t+w)+3p(t,w).
\]
Writing $a_j=\operatorname{advance}^j(0,0)$ for Cantor's diagonal
enumeration, the unpairing loop takes
\[
 U(n)=1+\sum_{j<n}r(a_j),\qquad
 r(a,b)=\begin{cases}3b+8,&a=0,\\5,&a>0.\end{cases}
\]
Its full restoration adds
$3\operatorname{hd}(n)+3\operatorname{tl}(n)+2$ steps.
The individual copy and times-four loops take $7v+2$ and $6v+1$
steps respectively.  Draining loops take $3v+1$ or $4v+1$.
The clear-only loop takes $2v+1$.  The gate comparison, on
$(t,h,x)$, has the primitive-recursive clock
\begin{align*}
 K_t(0,0)&=2t+3,& K_t(0,x+1)&=3x+5,\\
 K_t(h+1,0)&=2h+6,&K_t(h+1,x+1)&=3+K_t(h,x).
\end{align*}
It clears $h,x$ and preserves $t$ on inequality, but clears $t$
on equality.  This is a recursion through $\min(h,x)$: after
$j$ matching pairs, $h$ and $x$ have each decreased by $j$ and
exactly $3j$ rows have executed.  The two unmatched tails then
use the displayed draining clocks.  Gate node 142 takes $4a+1$
steps; the rest comprises $K_t(h,a)$, a $3a+1$ transfer from
$Y$ to $X$, a $3d+1$ transfer from $d$ to $h$, and the second
comparison.  Its intermediate clean-workspace conditions are
therefore accounted for explicitly.
Each formula comes from concatenating the displayed finite row
cycles.  Fix the input registers and induct on the number of
drained units or diagonal positions, expressing the current
registers and clock by their primitive-recursive iterates.  The
invariant is then quantifier-free.  This is the weak-base version
of the row-cycle arguments, including cases where the reference
Lean proof was written as induction generalizing a second register.
No universal quantifier over arbitrary environments is used as
an induction formula.

Fetch preserves $(t,w,d,h)$.  If $l=0$, its copy-and-restore prefix
takes $7n+4$ rows and leaves $(n,l)=(n+1,n)$, with the second $n$
meaning the original value; if $l>0$, it takes two rows and leaves
$(n,l)=(n,l-1)$.  It then unpairs this seed into $a,l$, with
$X=Y=0$, and uses one or two dispatch rows at 48--49.  This proves
exactly the scheduler seed and operand rules in (9).

The MP prelude reads $x$ into $d$, constructs $I(t,d+x)$ without
changing $w$, and drains this value into $X$.  The unchanged
comparison either reaches the success copy-$d$ block or the return
jump; it clears $t,h,X,Y$ on both paths.  Its equality case copies
$d+x$ to $t$ in $7(d+x)+2$ rows.  The only deletion is the dummy
push on that success path.  Thus its full contract is exactly the
new MP equation above for arbitrary $x,t,w,d,h$.
For its exact comparison clock use
\begin{align*}
 J(0,0)&=2,&J(0,x+1)&=3x+5,\\
 J(h+1,0)&=2h+6,&J(h+1,x+1)&=3+J(h,x).
\end{align*}
Unlike the gate comparison, equality exits immediately instead of
executing a clear-$t$ zero test.  Writing $b=d+x$ and $v=I(t,b)$,
the part after construction takes $3v+1+J(h,v)$ rows, then
$7b+2$ rows on success or zero on failure, and one return jump.

The AX core ends at 3519 with $t=G(H,D,c)$, $a=c$, $d=D$,
$h=H$, clean $X,Y$, and numerical stack
$[q_1,\ldots,q_9,t_{\rm old};w_{\rm old}]$.  Its two gate blocks
end at 3553 with $a=d=h=X=Y=0$ and $t=q_0$.  The last push
ends at 3572 with $t=0$ and the advertised eleven-entry stack.
The following MP block consumes the extra field $z$.  Since its
incoming $t=h=d=0$, it returns $t=0,h=0,d=z$ for both $z=0$
and $z>0$.  This proves the AX11 contract, including arbitrary
values of the fourth field, without assuming that field is zero.

The new Shift prefix at 3660 first clears $t$ in $2t+1$ rows,
then copies $h$ to $t$ in $7h+2$ rows while restoring $h$.
The final two rows set $X=1$ and jump to the unchanged comparison
at zero.  Its input is therefore $(t,h,X)=(h,h,1)$.  It halts
exactly when $h=1$, after $2t+7h+11$ rows in that case.
For $h\ne1$, it reaches the unpairing block at 11 after
$2t+9h+10$ rows, with $t=h_{\rm old}$, $h=X=Y=0$.
Unpairing $w$ then returns to 25 with the new Shift store.
These calculations also show that a command cannot halt inside
an arithmetic helper on the way to its specified endpoint.

Concatenate these traces at the displayed clean cutpoints.  All
segment clocks are fixed finite compositions and sums of the
primitive-recursive recurrences just given, evaluated on the
successive macro stores.  This gives a positive source segment
for every scheduler step and the separate terminating segment
for a halt.  The primary contracts are the same inherited
calculations with its original final AX return, MP dummy push
and sum guard.

For an explicit coarse majorant, an elementary routine with all
its numeric registers bounded by $B\geq1$ has duration and maximum
register bounded by $2^{20}(B+1)^4$, using the above recurrences.
One may iterate this majorant along the fixed finite macro plan,
including the parser and final guard routines.  Its iteration
count is the length of that materialized plan, not one application
per entire AX command.  In particular the nested axiom constructors
are not incorrectly assigned a degree-eight bound.  Summing the
actual recurrence costs is an alternative exact segment clock.

\subsection{Why the literal symbolic certificates are sound in iPRA}

The independent literal checker uses a two-sided finite-word tape
representation with arbitrary zero exterior.  A run atom represents
$0^a$ or $1^a$, where $a$ is an affine expression with nonnegative
coefficients in the gap, register values, and spare-list width.
The list blocks are
\[
                  B(v)=01^{v+1},\qquad C(v)=101^v.
\]
The four generic list traversals are derived from the literal rows
by first checking one arbitrary block, of length $v+2$, with both
external contexts left opaque.  Fix the code of a finite list $L$;
quantifier-free induction on the traversed prefix length concatenates
the recorded block traces, using the primitive-recursive functions
for the corresponding prefix and suffix of this fixed $L$.
This avoids induction on an unbounded assertion about all lists.
The only list-normalization
identity is
\[
                   C(L)1=1B(L),
\]
together with its reversal.  Its empty case is $1=1$, and its cons
case follows from $C(v)1=1B(v)$.  This is again an induction about
prefixes of a fixed primitive-recursively coded word, so is available
in iPRA.

There are precisely three operational symbolic rules.  A literal
row consumes one provably present bit.  A same-state,
same-direction self-loop traverses a homogeneous run of length
$a$; quantifier-free induction on $a$ gives its exact write effect,
duration and displacement.  A checked list traversal concatenates
the block lemma and has the recorded total width as its duration.
All other normalizations are word identities: merging equal-bit
runs, deleting an empty run, reversing a finite list, and absorbing
a finite zero suffix into the blank exterior.  None asserts a
runtime fact without a row or an induction just described.

The finite recorded derivations therefore prove the successor-frame
equation for every source site and both decrement branches, with
arbitrary nine natural registers, arbitrary nonnegative header
padding, and an arbitrary nonempty finite spare list.  The clocks
are the recorded affine expressions; the new table's 4,778 cases
all have positive constant part, at least three.  The frame
successors preserve nonnegative padding and a nonempty spare list;
increments append a spare block.  The fixed blank initialization
is a closed finite calculation of the entire nonzero tape support.
Finite row identities and closed startup calculations are
iPRA-provable without a uniform reflection principle for the checker.

It remains to reflect a physical halt rather than just simulate a
source step.  Augment the nine-register source configuration with
the header-padding counter, signed header/head location, and tagged
spare list used by its physical frame.  Update those auxiliary
fields by the finite
successor-frame contract.  They are initialized by the startup
receipt, and their nonnegativity and nonemptiness are the invariant
just established.  Forgetting the auxiliary fields commutes with
every source step.  Conversely, iterating their prescribed updates
lifts every source run from this initial configuration; no choice
or unbounded search is needed.  This augmentation therefore does
not change halting.  In the following, source configurations include
these fields, so that $F$ is a function of the current configuration,
not of an unspecified execution history.

Distinguish a live source configuration whose next
instruction is \texttt{HALT} from an executed-halt tag.  The former
still has its ordinary live control frame.  Executing that instruction
changes it to a tag containing the entire prehalt configuration;
the tag is thereafter absorbing.  Its physical realization is,
by definition, the endpoint of the recorded halt segment applied to
the preserved prehalt frame.  The finite halt contract proves that
this endpoint has the physical halt state; no prescribed final tape
shape is needed.  The endpoint is primitive recursive because the
segment clock and literal table iteration are primitive recursive.
Thus the halted-tag representation does not identify a prehalt
frame with a physical halt, and its halt correspondence is exact.

Let $F(c)$ be this extended full frame, $\ell(c)\geq1$ the
recorded segment clock for live $c$ (and one for an absorbing tag),
and $c_n$ the primitive source iteration.
Put $L_0=0$ and $L_{n+1}=L_n+\ell(c_n)$ while the source has
not executed its halt instruction, and continue with unit increments
afterwards.  The just-proved frame equations, including the
executed-halt tag equation, give
\[
 L_n\geq n,\qquad
 \operatorname{run}_M(b+L_n)=F(c_n),                       \tag{13}
\]
where $b$ is the fixed initialization time.  A source halt witness
therefore gives a physical halt witness.
Conversely, for a physical halt at time
$t$, bounded search among the first $t+1$ source configurations
either finds a source halt or gives, by (13), a nonhalt frame after
time $t$.  The latter contradicts absorption.  The search is
decidable and bounded, so the argument yields a source witness
intuitionistically; it uses no Markov principle.  Use the same
convention at the command-to-\emph{unaugmented}-source boundary:
a command halt first
executes its source segment and then stores that segment's exact
terminal source configuration in an absorbing tag.  It is not
identified with the live source entry at the halt-test command.
Apply the same positive-segment argument to this boundary, whose
frame is just the displayed nine-register clean command entry.
Then apply the deterministic ghost lift/forget equivalence, and
finally the augmented-source-to-physical argument (13).  In
particular, no physical padding, head position or spare list is
claimed to be recoverable from the scheduler store alone.

This proves each literal machine equivalent in iPRA to its
scheduler.  Equations (1) and (12) now prove the first assertion of
Theorem~\ref{thm:relative-ipra}; intuitionistic contraposition proves
the nonhalting assertion.

\subsection{The remaining absolute calibration}

No part of this proof translates an arbitrary internal PA proof
into a numerical-calculus refutation, or a numerical refutation
into a PA refutation.  The pinned repository does establish the
one-sided result $\Con(\PA)\to\NH(\text{modeled scheduler})$;
its README distinguishes that result from a physical-machine
calibration.  The missing direction needed for the full calibration
is, in particular, a verified internal transformation from arbitrary
PA refutations into numerical-calculus refutations.  The ordinary
PA calibration and the conditional iPRA endpoint do not supply
that transformation with its weak-base correctness proof.  Thus
the strengthened full equivalence requested by the user remains
open even though the relative theorem and the new literal
194-state construction are established here.


\section{What remains before the requested consistency theorem}
\label{sec:remaining}

At the pinned commit, the maintained PA README explicitly states that
neither physical-machine-specific direction of the desired iPRA
consistency equivalence is established for the primary PA194, original
PA195, or PA196 \cite{pastatus}. This is stronger than the observation that
some delivered packets lack a fresh Lean file. The repository identifies
specific open bridges: internal translation of arbitrary PA refutations
into the numerical calculus, and the all-model or internal physical
controller/source connection.

Substantial formal work does exist. The generic
\file{isigma1-machine} endpoint reduces the target to a machine-specific
calibration hypothesis, and \file{isigma1-transfer} supplies a route through
$\mathrm I\Sigma_1$ and conservation/negative translation \cite{isigmachine,isigtransfer}.
The modeled scheduler is calibrated to its target numerical derivations;
the target-to-PA direction has discharged its guarded separation premise.
These results concern explicitly identified intermediate objects and do
not silently identify them with a physical transition table.

The logical calibration can be stated without any implementation ambiguity.
For the exact primary numerical calculus $\mathcal C$ used here, one needs
primitive-recursive proof transformers $F,G$ whose following correctness
statements are themselves proved in iPRA:
\[
 \Prf_{\PA}(p,0=1)\ \Longrightarrow\ \Prf_{\mathcal C}(F(p),1),
 \qquad
 \Prf_{\mathcal C}(d,1)\ \Longrightarrow\ \Prf_{\PA}(G(d),0=1).
\]
The repository has already established the target-to-PA direction for its
modeled target. Reusing it here requires an explicit identification of the
proof predicates and encodings; the converse PA-to-target compiler remains
open. The conclusion codes must likewise be fixed.
A written PA correctness argument for a primitive-recursive function does
not by itself provide an iPRA correctness proof for the same function.
The unbounded-looking invariant ``for every environment there is a target
proof'' cannot simply be called quantifier-free induction. One promising
implementation is a bounded request tree recording every substitution and
proof-template instance needed by a particular input proof, with a single
decidable correctness invariant. This report does not complete that
compiler.

The absence of the finished repository theorem does not establish that the
requested iPRA statement is false. It prevents presenting that statement as
an inherited proved premise. Our fallback is consequently a concrete new
194-state design with explicit relative preservation, rather than an
unconditional completion of the user's strengthened target.

\section{Further attempts below 194 states}
\label{sec:search}

\subsection{Exact ordinary quotient calculation}

For each of the 82 spanning 194-state graphs, start with one block of
working controls and a separate terminal block. Repeatedly split working
controls by the ordered pair of (written symbol, direction, successor
block) signatures. The stable partition has 194 working classes for every
graph. This rules out strict action-preserving unconditional control
merging. It does not rule out a fold justified by tape context, a source
rewrite, or a different machine architecture.

\subsection{Permuting theorem emission and reassociating the decoder}

The source search allowed permutation of seven separately emitted theorem
heads as well as binary reassociation. Permutation changes selector ranks,
but not the available numerical axioms. Near the frontier, scoring checks
simultaneous fold compatibility rather than relying solely on an optimistic
maximum matching. The first recorded run made 3,631,115 proposals and
accepted 650,390 search moves in its 240-second budget. The additional run
used the H controller's provider transitions: 1,283,873 proposals and
221,857 accepted moves in 90 seconds. Both retained a best effective reader
count of 60. With 134 fixed controls, this is 194, not 193.

These two runs total 4,914,988 proposals. Their exact sources, seeds,
inputs, logs and counts are preserved. They are bounded heuristic searches,
not exhaustive analyses of all source programs or all theorem-head orders.

\subsection{Local physical synthesis on the new composition}

A corrected source-context collector includes the interiors of accelerated
constant-symbol scans. On PA194AX11QH it obtains 8,462 radius-four windows,
with 6,242 distinct accelerated-scan observations, while replaying all
4,778 source outcomes. Arithmetic-box interiors are not narrowed by the
source-only collector. Startup is intentionally excluded at this stage.

The two-row return synthesizer considers all $\binom{194}{2}=18,721$
control pairs, allowing repeated visits to the merged control. With the
fixed radius-four envelope and a 120-transition closure bound, it makes
23,902,664 partial action trials and checks 143,951 completed action
combinations. No quotient survives. This is a result for that local
synthesis family and its envelope, not a general nonexistence theorem.

A second screen first enumerates 357 one-row detours that rejoin an old
transition endpoint on the sampled symbolic envelope, then reruns the
quotient synthesis after each proposed detour. All 357 screens finish
with zero hits. The second stage deliberately reuses the original context
envelope after its first edit; consequently it is a heuristic screen.
Any positive result would have needed a fresh complete source proof and
blank initialization. A negative result from this restricted procedure
cannot exclude other simultaneous edits or longer simulations.

\subsection{The inherited 193-state marker reduction on the new table}

Apply the familiar native-label edit to the new composition:
identify 301 with 88, redirect targets accordingly, delete 301, and set
\[
 (88,0)\mapsto(0,L,51),\qquad(51,0)\mapsto(0,R,122).
\]
This yields a 193-control experimental table. Choosing any of its 193
controls as a blank initial state produces two halting computations and
191 certified outward rays, all classified within 93 transitions. At the
usual start 74 the machine reaches state 35 after 92 transitions, at head
14 with finite one-support $\{-3,1\}$. The row $35,0\mapsto1R35$ then
writes ones forever into the untouched right-hand exterior. This is an
exact infinite-continuation certificate, not a fuel-exhaustion verdict.
It cannot repair the consistency target.

The zero saving found by these searches is reported candidly. The record
194 is not proved minimal, and the complete search space is vastly larger
than the families examined here.

\section{Lessons from the other suggested projects}

The ZFC work is useful for contextual history folds, independently derived
register-block lemmas, exact whole-tape startup, and persistent-store
fairness. Its reviewed marker-relay attempts also warn against transferring
a locally successful rewrite without checking every source outcome.
Its consistency target is different: a Con(ZFC) calibration cannot be
substituted for a Con(PA) calibration \cite{zfc}.

The DiophantineConsistency project supplies the relevant proof-strength
methodology: pin the exact arithmetic sentence, give finite witness
transformations, and distinguish standard-model equivalence from native
iPRA derivability \cite{dio}. A compact polynomial, few variables, or low
arithmetic-operation count is not a count of binary Turing controls;
constants, enumeration, decoding and blank initialization still have to
be implemented. No state saving from this route is asserted here.

ConciseIPRA supplies compact presentations of the \emph{metatheory},
including methods that derive logic and induction interfaces from smaller
primitives \cite{concise}. Such interfaces may help organize an eventual
proof checker or compiler. Replacing the PA refutation search with an iPRA
refutation search changes the consistency statement being encoded. A
smaller checker for the latter does not solve the present problem.

\section{Package, reproduction and trust}

All primary replay commands use Python 3.10 or newer and its standard
library. Assertions must remain enabled. The local reduction and grammar
searches additionally need a C++17 compiler. No search is needed to
reconstruct the delivered candidate. No network access is needed for its
verification, the old-table replay, the frozen-corpus novelty check, or
the finite premise-exchange certificate.

From the package root:
\begin{lstlisting}
python3 -B verify_all.py --output fresh_checks
\end{lstlisting}
The driver reruns the complete new construction, command/source checks,
independent literal source contracts, old primary replay, exact corpus
novelty, finite restricted-premise exchange, and startup-ray certificates.
Its receipt explicitly states that these tests do not certify the missing
Con(PA) calibration. The complete corpus census can also be rerun from a
checkout of the pinned commit; instructions and the recognition scope are
in \file{comparison/README.txt}.

To rebuild this article:
\begin{lstlisting}
cd report
python3 generate_tables.py
pdflatex -interaction=nonstopmode -halt-on-error report.tex
pdflatex -interaction=nonstopmode -halt-on-error report.tex
\end{lstlisting}
The package retains imported licenses and a provenance manifest. The
transition tables are exact bytes, and \file{SHA256SUMS} identifies all
shipped files except itself. Historical machine names and proof claims
remain attributed to their source packets. No proof-assistant endpoint is
invented to fill an unfinished logical bridge.

\appendix
\section{Every principal 194-state design}
\label{app:principal}
All rows below have 194 working states. $r$ is the numerical constructor
coefficient; Q marks premise exchange; the diagonal column records its
parameter; AX is prefix length. MP is P (visible result with dummy push),
C (cached result), or N (visible result without push). Halt is the sum or
cache test. Sym. marks the swapped comparison order. The detailed
mechanisms and full identifiers are in \file{principal_designs40.csv}.
Every historical iPRA status remains ``not established for the literal
physical machine.''
\small
\begin{longtable}{@{}p{4.15cm}ccccccc@{}}
\toprule
Design & $r$ & Q & $\Gamma$ & AX & MP & Halt & Sym.\\
\midrule\endfirsthead
\toprule Design & $r$ & Q & $\Gamma$ & AX & MP & Halt & Sym.\\\midrule\endhead
\bottomrule\endfoot
M4R & 4 & -- & A & 10 & P & sum & --\\
M5 & 5 & -- & A & 10 & P & sum & --\\
M6 & 6 & -- & A & 10 & P & sum & --\\
PA194 (68/8) & 4 & -- & A & 10 & P & sum & --\\
PA194-AD & 4 & -- & A & 10 & P & sum & --\\
PA194A & 4 & -- & A & 10 & C & cache & yes\\
PA194AX11 & 4 & -- & A & 11 & N & sum & --\\
PA194C & 4 & -- & A & 10 & C & cache & --\\
PA194CS & 4 & -- & A & 10 & C & cache & yes\\
PA194CSZ = PA194X & 4 & -- & A & 10 & C & cache & yes\\
PA194D & 4 & -- & D & 10 & C & cache & --\\
PA194DHS & 4 & -- & D & 10 & N & cache & yes\\
PA194F & 4 & -- & A & 10 & C & cache & yes\\
PA194H & 4 & -- & A & 10 & N & cache & --\\
PA194H5 & 5 & -- & A & 10 & N & cache & --\\
PA194H6 & 6 & -- & A & 10 & N & cache & --\\
PA194HA & 4 & -- & A & 10 & C & sum & --\\
PA194HAZ & 4 & -- & A & 10 & C & sum & --\\
PA194HS = filed PA194NS & 4 & -- & A & 10 & N & cache & yes\\
PA194J & 4 & -- & A & 10 & C & cache & yes\\
PA194N & 4 & -- & A & 10 & N & sum & --\\
PA194N5 & 5 & -- & A & 10 & N & sum & --\\
PA194N6 & 6 & -- & A & 10 & N & sum & --\\
PA194NS (sum guard) & 4 & -- & A & 10 & N & sum & yes\\
PA194Q & 4 & yes & A & 10 & C & cache & --\\
PA194Q5 & 5 & yes & A & 10 & P & sum & --\\
PA194Q6 & 6 & yes & A & 10 & P & sum & --\\
PA194QH & 4 & yes & A & 10 & N & cache & --\\
PA194QH5 & 5 & yes & A & 10 & N & cache & --\\
PA194QH6 & 6 & yes & A & 10 & N & cache & --\\
PA194QHS & 4 & yes & A & 10 & N & cache & yes\\
PA194QN = PA194NQ & 4 & yes & A & 10 & N & sum & --\\
PA194QN5 & 5 & yes & A & 10 & N & sum & --\\
PA194QN6 & 6 & yes & A & 10 & N & sum & --\\
PA194QNS & 4 & yes & A & 10 & N & sum & yes\\
PA194QS & 4 & yes & A & 10 & P & sum & yes\\
PA194R & 4 & -- & A & 10 & C & cache & yes\\
PA194SUMS & 4 & -- & A & 10 & P & sum & yes\\
PA194ZR & 4 & -- & A & 10 & C & cache & yes\\
Primary PA194 (67/7) & 4 & -- & A & 10 & P & sum & --\\
\midrule\textbf{New AX11QH} & 4 & yes & A & 11 & N & cache & --\\
\end{longtable}
\normalsize


\section{The 82 spanning frontier graphs}
\label{app:frontier}
The following identifiers are prefixes of full strict canonical hashes,
not hashes of arbitrary table bytes. The complete tuples, full hashes,
source locators and evidence annotations are retained in the comparison
artifacts. P: reported principal; S: supplementary; V: source-validated
rotation; A: research auxiliary; R: explicitly rejected. These are historical
categories, not new consistency proofs. The new machine is additional to
this pinned corpus.
\small\begin{longtable}{@{}p{2.3cm}p{7.8cm}p{2.25cm}@{}}
\toprule Graph prefix & Reported name or role & Evidence\\\midrule\endfirsthead
\toprule Graph prefix & Reported name or role & Evidence\\\midrule\endhead
\bottomrule\endfoot
\texttt{a375028ad457} & Rejected append-zero tail & R\\
\texttt{ff86e0ca1cee} & Rejected comparator return 103 & R\\
\texttt{feea621714b0} & Rejected dec6 & R\\
\texttt{52ca5314cd73} & Rejected node1-X & R\\
\texttt{f0b9edcb257d} & Rejected short scanner & R\\
\texttt{a3dfc75acd0d} & M4R & P\\
\texttt{b7c97ee64c25} & M5 & P\\
\texttt{0db332b9e9c5} & M6 & P\\
\texttt{b98fffb6310a} & PA194 (68/8) & P\\
\texttt{f5efe9523f57} & PA194-AD & P\\
\texttt{01aa08f9cf90} & PA194A & P\\
\texttt{1c8542f6445e} & PA194AX11 & P\\
\texttt{da745794648f} & PA194C & P\\
\texttt{7f0f46bef40e} & PA194CS & P\\
\texttt{15959ba90310} & PA194CSZ = PA194X & P\\
\texttt{462b00a5ac99} & PA194D & P\\
\texttt{2b2be6b73ba2} & PA194DHS & P\\
\texttt{b8f45d21e029} & PA194F & P\\
\texttt{46becfb1d977} & PA194H & P\\
\texttt{984c271bcb9b} & PA194H5 & P\\
\texttt{960f5ae293dd} & PA194H6 & P\\
\texttt{b3be4de9519d} & PA194HA & P\\
\texttt{a36565c03297} & PA194HAZ & P\\
\texttt{fd0f6a53ee42} & PA194HS = filed PA194NS & P\\
\texttt{ba88eb1be4ae} & PA194J & P\\
\texttt{5f4126e5adc2} & PA194N & P\\
\texttt{6a5f3270f085} & PA194N5 & P\\
\texttt{89fc11dc5ced} & PA194N6 & P\\
\texttt{701fe9c7ac14} & PA194NS (sum guard) & P\\
\texttt{8fee94fb004d} & PA194Q & P\\
\texttt{d36874f3de2f} & PA194Q5 & P\\
\texttt{1040e3d8e079} & PA194Q6 & P\\
\texttt{6bef585d5c78} & PA194QH & P\\
\texttt{ec1d05c85a97} & PA194QH5 & P\\
\texttt{90f17156c1ea} & PA194QH6 & P\\
\texttt{1ecac647dcc5} & PA194QHS & P\\
\texttt{349902453837} & PA194QN = PA194NQ & P\\
\texttt{4a574f192796} & PA194QN5 & P\\
\texttt{d950adad01c5} & PA194QN6 & P\\
\texttt{0db6cfbd079a} & PA194QNS & P\\
\texttt{9c0f753a5f71} & PA194QS & P\\
\texttt{7d7b3a970c8d} & PA194R & P\\
\texttt{b802032761b6} & PA194SUMS & P\\
\texttt{307b9f6064be} & PA194ZR & P\\
\texttt{0b1f3bd2a6aa} & Primary PA194 (67/7) & P\\
\texttt{e9205ed8d7d7} & C-I-Z & A\\
\texttt{347ec209e98e} & CS both-zero6 & A\\
\texttt{3969780d1528} & MPH auxiliary & A\\
\texttt{4078bc70be6b} & N6 append-zero exploratory patch & A\\
\texttt{ba8e4badf4ff} & PA194-altC & A\\
\texttt{c87f12477e42} & PA194N-AD (auxiliary) & A\\
\texttt{843758c26679} & Q with original controller (QS parent) & A\\
\texttt{0752d92b9b37} & Source rotation 100\_164\_left & A\\
\texttt{20ee651245a6} & Source rotation 166\_179\_left & A\\
\texttt{16731549275d} & Source rotation 19\_29\_left & A\\
\texttt{e10fd4a65532} & Source rotation 2000\_2001\_right & A\\
\texttt{4495c1132037} & Source rotation 2007\_2008\_left & A\\
\texttt{543c6fbe4fbe} & Source rotation 2009\_2011\_left & A\\
\texttt{faa3eda5bc88} & Source rotation 2040\_2041\_right & A\\
\texttt{cfe7c2558157} & Source rotation 2041\_2042\_right & A\\
\texttt{7167be4a487a} & Source rotation 2041\_2047\_left & A\\
\texttt{488b15a851c0} & Source rotation 2047\_2048\_right & A\\
\texttt{13092ab9697f} & Source rotation 2047\_2053\_left & A\\
\texttt{9e841694cec6} & Source rotation 2048\_2049\_left & A\\
\texttt{e9f165e9cd5c} & Source rotation 2053\_2054\_right & A\\
\texttt{b3e21774a966} & Source rotation 214\_312\_left & A\\
\texttt{2b8bfc9bb21a} & Source rotation 217\_197\_right & A\\
\texttt{e9584acb96eb} & Source rotation 221\_217\_left & A\\
\texttt{d98716205de4} & Source rotation 29\_44\_left & A\\
\texttt{137c822f17cd} & Source rotation 312\_38\_left & A\\
\texttt{b58cd23e9c63} & Source rotation 100\_174\_right & V\\
\texttt{eb2d614d50ea} & Source rotation 2008\_2009\_right & V\\
\texttt{4332ea11fbbf} & Source rotation 2009\_2010\_right & V\\
\texttt{58b6850e764e} & Source rotation 2011\_2012\_left & V\\
\texttt{216161663948} & Source rotation 2061\_2062\_left & V\\
\texttt{68c5e616f092} & CSZA final-register 3 & S\\
\texttt{66bd66525394} & CSZA final-register 5 & S\\
\texttt{1b63db82fcb7} & CSZA final-register 6 & S\\
\texttt{11a205ab00b0} & Gamma(D) + CSZA & S\\
\texttt{54b876cc9ebb} & PA194CZ & S\\
\texttt{38aed066c1c4} & PA194F\_AD & S\\
\texttt{828e45657d1c} & PA194R\_AD & S\\
\end{longtable}\normalsize


\section{Complete transition table of PA194AX11QH}
\label{app:table}
Initial state is $0$ and the tape is all zero. In an entry $wDq'$, write
$w$, move in direction $D$, and enter $q'$. The terminal symbol $H$ is label
194 and has no outgoing rows. Both panels belong to one 194-state machine.
\small\begin{longtable}{@{}rcc@{\hspace{1.6cm}}rcc@{}}
\toprule $q$ & Read 0 & Read 1 & $q$ & Read 0 & Read 1\\\midrule\endfirsthead
\toprule $q$ & Read 0 & Read 1 & $q$ & Read 0 & Read 1\\\midrule\endhead
\bottomrule\endfoot
0 & $1\mathrm{R}1$ & $1\mathrm{R}2$ & 97 & $0\mathrm{R}36$ & $1\mathrm{R}112$\\
1 & $1\mathrm{R}3$ & $1\mathrm{R}1$ & 98 & $0\mathrm{R}111$ & $1\mathrm{R}75$\\
2 & $0\mathrm{R}0$ & $1\mathrm{R}2$ & 99 & $0\mathrm{L}63$ & $0\mathrm{R}109$\\
3 & $0\mathrm{R}4$ & $0\mathrm{R}1$ & 100 & $0\mathrm{L}62$ & $1\mathrm{R}113$\\
4 & $1\mathrm{R}5$ & $1\mathrm{L}6$ & 101 & $0\mathrm{L}90$ & $1\mathrm{R}114$\\
5 & $1\mathrm{L}7$ & $1\mathrm{L}8$ & 102 & $0\mathrm{R}115$ & $0\mathrm{L}14$\\
6 & $0\mathrm{L}7$ & $1\mathrm{L}6$ & 103 & $0\mathrm{R}69$ & $1\mathrm{R}103$\\
7 & $0\mathrm{L}9$ & $1\mathrm{L}6$ & 104 & $0\mathrm{R}116$ & $1\mathrm{R}95$\\
8 & $0\mathrm{L}10$ & $0\mathrm{L}11$ & 105 & $0\mathrm{R}116$ & $1\mathrm{R}67$\\
9 & $1\mathrm{R}12$ & $0\mathrm{L}9$ & 106 & $0\mathrm{R}117$ & $0\mathrm{L}14$\\
10 & $0\mathrm{L}9$ & $0\mathrm{L}10$ & 107 & $0\mathrm{R}118$ & $1\mathrm{L}71$\\
11 & $1\mathrm{L}8$ & $1\mathrm{L}11$ & 108 & $0\mathrm{R}119$ & $1\mathrm{R}120$\\
12 & $0\mathrm{R}12$ & $0\mathrm{R}13$ & 109 & $0\mathrm{R}14$ & $1\mathrm{R}109$\\
13 & $0\mathrm{L}14$ & $0\mathrm{R}13$ & 110 & $0\mathrm{R}66$ & $1\mathrm{R}110$\\
14 & $1\mathrm{L}15$ & $0\mathrm{L}15$ & 111 & $0\mathrm{R}121$ & $1\mathrm{L}35$\\
15 & $0\mathrm{L}15$ & $1\mathrm{L}16$ & 112 & $0\mathrm{R}122$ & $1\mathrm{R}123$\\
16 & $0\mathrm{L}17$ & $1\mathrm{L}16$ & 113 & $0\mathrm{R}88$ & $0\mathrm{L}14$\\
17 & $0\mathrm{L}18$ & $1\mathrm{L}16$ & 114 & $0\mathrm{R}90$ & $1\mathrm{L}95$\\
18 & $0\mathrm{L}19$ & $1\mathrm{L}16$ & 115 & $0\mathrm{R}124$ & $1\mathrm{R}124$\\
19 & $0\mathrm{L}20$ & $1\mathrm{L}16$ & 116 & $0\mathrm{L}63$ & $0\mathrm{R}103$\\
20 & $0\mathrm{L}21$ & $1\mathrm{L}16$ & 117 & $0\mathrm{R}68$ & $0\mathrm{R}91$\\
21 & $0\mathrm{L}22$ & $1\mathrm{L}16$ & 118 & $0\mathrm{R}58$ & $0\mathrm{R}91$\\
22 & $0\mathrm{L}23$ & $1\mathrm{L}16$ & 119 & $0\mathrm{L}86$ & $1\mathrm{R}125$\\
23 & $0\mathrm{R}23$ & $1\mathrm{R}24$ & 120 & $0\mathrm{R}126$ & $1\mathrm{R}127$\\
24 & $0\mathrm{R}25$ & $1\mathrm{R}26$ & 121 & $0\mathrm{R}58$ & $0\mathrm{R}110$\\
25 & $0\mathrm{R}27$ & $1\mathrm{R}28$ & 122 & $0\mathrm{R}99$ & $1\mathrm{R}113$\\
26 & $0\mathrm{R}29$ & $1\mathrm{R}30$ & 123 & $0\mathrm{R}79$ & $1\mathrm{R}78$\\
27 & $0\mathrm{R}31$ & $1\mathrm{R}H$ & 124 & $0\mathrm{R}88$ & $1\mathrm{R}88$\\
28 & $0\mathrm{L}32$ & $1\mathrm{R}33$ & 125 & $0\mathrm{R}128$ & $1\mathrm{R}129$\\
29 & $0\mathrm{R}34$ & $1\mathrm{R}35$ & 126 & $0\mathrm{R}32$ & $1\mathrm{R}36$\\
30 & $0\mathrm{L}36$ & $1\mathrm{L}37$ & 127 & $0\mathrm{R}31$ & $1\mathrm{R}98$\\
31 & $0\mathrm{R}38$ & $1\mathrm{R}39$ & 128 & $0\mathrm{R}130$ & $1\mathrm{R}131$\\
32 & $0\mathrm{R}40$ & $1\mathrm{R}41$ & 129 & $0\mathrm{R}132$ & $1\mathrm{R}133$\\
33 & $0\mathrm{R}42$ & $1\mathrm{R}43$ & 130 & $0\mathrm{R}134$ & $1\mathrm{R}135$\\
34 & $0\mathrm{R}44$ & $1\mathrm{R}45$ & 131 & $0\mathrm{R}136$ & $1\mathrm{R}137$\\
35 & $0\mathrm{R}46$ & $1\mathrm{R}33$ & 132 & $0\mathrm{R}138$ & $1\mathrm{R}107$\\
36 & $0\mathrm{R}47$ & $1\mathrm{R}48$ & 133 & $0\mathrm{R}139$ & $1\mathrm{R}97$\\
37 & $0\mathrm{R}49$ & $1\mathrm{R}50$ & 134 & $0\mathrm{R}97$ & $1\mathrm{R}35$\\
38 & $0\mathrm{R}51$ & $1\mathrm{R}42$ & 135 & $0\mathrm{R}140$ & $1\mathrm{R}124$\\
39 & $0\mathrm{R}52$ & $1\mathrm{R}53$ & 136 & $0\mathrm{R}141$ & $1\mathrm{R}142$\\
40 & $0\mathrm{R}54$ & $1\mathrm{R}55$ & 137 & $0\mathrm{R}143$ & $1\mathrm{R}144$\\
41 & $0\mathrm{R}56$ & $1\mathrm{R}57$ & 138 & $0\mathrm{R}145$ & $1\mathrm{R}41$\\
42 & $0\mathrm{R}58$ & $0\mathrm{R}59$ & 139 & $0\mathrm{R}146$ & $1\mathrm{R}147$\\
43 & $0\mathrm{R}60$ & $0\mathrm{L}14$ & 140 & $0\mathrm{L}148$ & $1\mathrm{L}143$\\
44 & $0\mathrm{R}61$ & $1\mathrm{R}62$ & 141 & $0\mathrm{R}28$ & $1\mathrm{R}149$\\
45 & $0\mathrm{R}63$ & $1\mathrm{R}52$ & 142 & $0\mathrm{R}135$ & $1\mathrm{R}150$\\
46 & $0\mathrm{R}64$ & $1\mathrm{R}65$ & 143 & $0\mathrm{R}151$ & $1\mathrm{R}152$\\
47 & $0\mathrm{R}58$ & $0\mathrm{R}66$ & 144 & $0\mathrm{R}153$ & $1\mathrm{R}154$\\
48 & $0\mathrm{R}67$ & $0\mathrm{L}14$ & 145 & $0\mathrm{L}130$ & $1\mathrm{R}88$\\
49 & $0\mathrm{R}68$ & $0\mathrm{R}69$ & 146 & $0\mathrm{R}31$ & $1\mathrm{R}52$\\
50 & $0\mathrm{R}70$ & $0\mathrm{L}48$ & 147 & $0\mathrm{R}155$ & $1\mathrm{R}146$\\
51 & $0\mathrm{R}64$ & $1\mathrm{R}53$ & 148 & $0\mathrm{R}156$ & $1\mathrm{R}157$\\
52 & $0\mathrm{L}71$ & $0\mathrm{L}15$ & 149 & $0\mathrm{R}158$ & $1\mathrm{R}159$\\
53 & $0\mathrm{R}72$ & $0\mathrm{L}10$ & 150 & $0\mathrm{L}159$ & $1\mathrm{R}160$\\
54 & $0\mathrm{R}73$ & $1\mathrm{R}74$ & 151 & $0\mathrm{R}41$ & $1\mathrm{L}161$\\
55 & $0\mathrm{R}75$ & $1\mathrm{R}76$ & 152 & $0\mathrm{R}162$ & $1\mathrm{R}163$\\
56 & $0\mathrm{R}77$ & $1\mathrm{R}78$ & 153 & $0\mathrm{R}164$ & $1\mathrm{R}164$\\
57 & $0\mathrm{R}79$ & $1\mathrm{R}80$ & 154 & $0\mathrm{R}165$ & $1\mathrm{R}166$\\
58 & $0\mathrm{L}81$ & $0\mathrm{R}1$ & 155 & $0\mathrm{R}92$ & $1\mathrm{R}111$\\
59 & $0\mathrm{R}82$ & $1\mathrm{R}59$ & 156 & $0\mathrm{R}76$ & $1\mathrm{R}165$\\
60 & $0\mathrm{L}73$ & $1\mathrm{L}46$ & 157 & $0\mathrm{R}167$ & $1\mathrm{R}165$\\
61 & $0\mathrm{R}83$ & $1\mathrm{R}84$ & 158 & $0\mathrm{R}168$ & $1\mathrm{R}169$\\
62 & $0\mathrm{R}54$ & $1\mathrm{R}85$ & 159 & $0\mathrm{R}170$ & $1\mathrm{R}157$\\
63 & $0\mathrm{L}86$ & $1\mathrm{R}87$ & 160 & $0\mathrm{R}135$ & $1\mathrm{R}130$\\
64 & $0\mathrm{R}58$ & $0\mathrm{R}82$ & 161 & $0\mathrm{R}138$ & $1\mathrm{R}135$\\
65 & $0\mathrm{R}88$ & $1\mathrm{R}30$ & 162 & $0\mathrm{R}171$ & $1\mathrm{R}172$\\
66 & $0\mathrm{R}89$ & $1\mathrm{R}66$ & 163 & $0\mathrm{R}173$ & $1\mathrm{R}172$\\
67 & $0\mathrm{L}90$ & $1\mathrm{L}33$ & 164 & $0\mathrm{L}174$ & $1\mathrm{R}175$\\
68 & $0\mathrm{L}81$ & $0\mathrm{R}2$ & 165 & $0\mathrm{L}162$ & $1\mathrm{L}176$\\
69 & $0\mathrm{R}91$ & $1\mathrm{R}69$ & 166 & $0\mathrm{R}161$ & $1\mathrm{R}176$\\
70 & $0\mathrm{R}72$ & $1\mathrm{R}92$ & 167 & $0\mathrm{R}97$ & $1\mathrm{R}177$\\
71 & $0\mathrm{R}46$ & $1\mathrm{R}93$ & 168 & $0\mathrm{R}134$ & $1\mathrm{R}178$\\
72 & $0\mathrm{R}42$ & $0\mathrm{L}13$ & 169 & $0\mathrm{R}150$ & $1\mathrm{R}163$\\
73 & $0\mathrm{R}94$ & $1\mathrm{R}37$ & 170 & $0\mathrm{R}179$ & $1\mathrm{R}180$\\
74 & $0\mathrm{R}30$ & $1\mathrm{R}95$ & 171 & $0\mathrm{R}167$ & $1\mathrm{R}140$\\
75 & $0\mathrm{R}42$ & $1\mathrm{R}96$ & 172 & $0\mathrm{R}134$ & $1\mathrm{R}145$\\
76 & $0\mathrm{R}97$ & $1\mathrm{R}98$ & 173 & $0\mathrm{R}181$ & $1\mathrm{R}41$\\
77 & $0\mathrm{R}99$ & $1\mathrm{L}100$ & 174 & $0\mathrm{R}175$ & $1\mathrm{R}182$\\
78 & $0\mathrm{R}30$ & $1\mathrm{R}77$ & 175 & $0\mathrm{R}167$ & $1\mathrm{R}145$\\
79 & $0\mathrm{R}93$ & $1\mathrm{R}101$ & 176 & $0\mathrm{R}183$ & $1\mathrm{R}173$\\
80 & $0\mathrm{R}30$ & $1\mathrm{R}102$ & 177 & $0\mathrm{R}107$ & $1\mathrm{R}85$\\
81 & $1\mathrm{R}81$ & $1\mathrm{L}15$ & 178 & $0\mathrm{R}184$ & $1\mathrm{R}165$\\
82 & $0\mathrm{R}103$ & $1\mathrm{R}82$ & 179 & $0\mathrm{R}134$ & $1\mathrm{R}181$\\
83 & $0\mathrm{R}94$ & $0\mathrm{L}7$ & 180 & $0\mathrm{R}185$ & $1\mathrm{R}185$\\
84 & $0\mathrm{R}104$ & $1\mathrm{R}105$ & 181 & $0\mathrm{R}167$ & $1\mathrm{R}151$\\
85 & $0\mathrm{R}42$ & $1\mathrm{R}106$ & 182 & $0\mathrm{R}151$ & $1\mathrm{R}130$\\
86 & $0\mathrm{R}107$ & $1\mathrm{R}48$ & 183 & $0\mathrm{R}164$ & $1\mathrm{R}186$\\
87 & $0\mathrm{R}108$ & $0\mathrm{L}72$ & 184 & $0\mathrm{R}140$ & $1\mathrm{R}187$\\
88 & $0\mathrm{R}68$ & $0\mathrm{R}66$ & 185 & $0\mathrm{R}188$ & $1\mathrm{L}189$\\
89 & $0\mathrm{R}109$ & $1\mathrm{R}89$ & 186 & $0\mathrm{R}190$ & $1\mathrm{R}41$\\
90 & $0\mathrm{R}49$ & $1\mathrm{R}100$ & 187 & $0\mathrm{R}32$ & $1\mathrm{R}189$\\
91 & $0\mathrm{R}110$ & $1\mathrm{R}91$ & 188 & $0\mathrm{R}88$ & $1\mathrm{R}191$\\
92 & $0\mathrm{R}30$ & $1\mathrm{R}48$ & 189 & $0\mathrm{R}148$ & $1\mathrm{R}192$\\
93 & $0\mathrm{R}47$ & $1\mathrm{R}60$ & 190 & $0\mathrm{R}41$ & $1\mathrm{R}138$\\
94 & $0\mathrm{L}63$ & $0\mathrm{R}89$ & 191 & $0\mathrm{R}174$ & $1\mathrm{R}88$\\
95 & $0\mathrm{R}94$ & $1\mathrm{L}111$ & 192 & $0\mathrm{R}182$ & $1\mathrm{R}193$\\
96 & $0\mathrm{R}101$ & $0\mathrm{L}14$ & 193 & $0\mathrm{R}161$ & $1\mathrm{R}138$\\
\end{longtable}\normalsize


\begin{thebibliography}{99}
\bibitem{repo} Vladimir Reshetnikov et al., \emph{ProveIt}, pinned commit
\snapshot, \url{https://github.com/VladimirReshetnikov/ProveIt/tree/1c9524d8bf2e0a27b9a16ff622359927f536e663}.
\bibitem{pastatus} \emph{PA consistency machines: current proof status},
\url{https://github.com/VladimirReshetnikov/ProveIt/blob/1c9524d8bf2e0a27b9a16ff622359927f536e663/Computability/BusyBeaver/Research/PA/README.md}.
\bibitem{packets} \emph{Delivered research packets on Con(PA) machines},
\url{https://github.com/VladimirReshetnikov/ProveIt/blob/1c9524d8bf2e0a27b9a16ff622359927f536e663/Computability/BusyBeaver/Research/PA/docs/README.md}.
\bibitem{machineledger} \emph{Machines in the delivered packets},
\url{https://github.com/VladimirReshetnikov/ProveIt/blob/1c9524d8bf2e0a27b9a16ff622359927f536e663/Computability/BusyBeaver/Research/PA/docs/MACHINES.md}.
\bibitem{primary} \emph{Primary PA194 research packet},
\url{https://github.com/VladimirReshetnikov/ProveIt/tree/1c9524d8bf2e0a27b9a16ff622359927f536e663/Computability/BusyBeaver/Research/PA/docs/pa194_conpa_2026-11-13}.
\bibitem{incoming} \emph{Incoming research archives},
\url{https://github.com/VladimirReshetnikov/ProveIt/tree/1c9524d8bf2e0a27b9a16ff622359927f536e663/docs/incoming}.
\bibitem{isigmachine} \emph{Generic machine calibration and iPRA endpoint},
\url{https://github.com/VladimirReshetnikov/ProveIt/tree/1c9524d8bf2e0a27b9a16ff622359927f536e663/Computability/BusyBeaver/Research/PA/isigma1-machine}.
\bibitem{isigtransfer} \emph{Transfer from I\,$\Sigma_1$},
\url{https://github.com/VladimirReshetnikov/ProveIt/tree/1c9524d8bf2e0a27b9a16ff622359927f536e663/Computability/BusyBeaver/Research/PA/isigma1-transfer}.
\bibitem{zfc} \emph{Smaller machines for Con(ZFC)},
\url{https://github.com/VladimirReshetnikov/ProveIt/tree/1c9524d8bf2e0a27b9a16ff622359927f536e663/Computability/BusyBeaver/Research/ZFC}.
\bibitem{dio} \emph{DiophantineConsistency},
\url{https://github.com/VladimirReshetnikov/ProveIt/tree/1c9524d8bf2e0a27b9a16ff622359927f536e663/Computability/DiophantineConsistency}.
\bibitem{concise} \emph{Concise presentations of intuitionistic primitive recursive arithmetic},
\url{https://github.com/VladimirReshetnikov/ProveIt/tree/1c9524d8bf2e0a27b9a16ff622359927f536e663/Logic/PrimitiveRecursiveArithmetic/ConciseIPRA}.
\end{thebibliography}
\end{document}
